% =====================================================================
% Phase 2 — RAG sur la documentation Python officielle
% Concepts essentiels du Retrieval-Augmented Generation avec LangChain.
% Compilation : pdflatex -shell-escape phase2.tex (deux fois pour la TOC)
% Pré-requis  : pygmentize (pip install Pygments) + paquet TeX "minted"
% =====================================================================
\documentclass[11pt,a4paper]{article}

% ---------- Encodage & langue ----------
\usepackage[utf8]{inputenc}
\usepackage[T1]{fontenc}
\usepackage[french]{babel}
\usepackage{lmodern}
\usepackage{microtype}

% Glyphes Unicode utilisés dans le code Python copié-collé
\DeclareUnicodeCharacter{2713}{\textbf{V}}   % ✓
\DeclareUnicodeCharacter{2717}{\textbf{X}}   % ✗
\DeclareUnicodeCharacter{2192}{$\rightarrow$}% →
\DeclareUnicodeCharacter{2500}{-}            % ─ box drawing
\DeclareUnicodeCharacter{276F}{>}            % ❯

% ---------- Mise en page ----------
\usepackage[a4paper,margin=2.2cm]{geometry}
\usepackage{parskip}
\usepackage{xcolor}
\usepackage{enumitem}
\usepackage{booktabs}
\usepackage{hyperref}
\hypersetup{
  colorlinks=true,
  linkcolor=blue!60!black,
  urlcolor=blue!60!black,
  citecolor=blue!60!black,
  bookmarksopen=true,
  bookmarksopenlevel=2,
}

% ---------- Boîtes & callouts ----------
\usepackage[most]{tcolorbox}
\tcbset{
  enhanced,
  colback=blue!3,
  colframe=blue!50!black,
  fonttitle=\bfseries,
  arc=2pt, boxrule=0.6pt,
}

% ---------- Code avec minted ----------
\usepackage{minted}
\setminted{
  fontsize=\small,
  linenos=true,
  numbersep=6pt,
  frame=lines,
  framesep=2mm,
  bgcolor=gray!4,
  breaklines=true,
  breakanywhere=true,
  python3=true,
  tabsize=4,
}
\newminted{python}{}

% ---------- TikZ ----------
\usepackage{tikz}
\usetikzlibrary{
  arrows.meta, positioning, shapes.geometric, shapes.misc,
  fit, backgrounds, calc, trees,
  decorations.pathreplacing, decorations.markings,
}

\tikzset{
  node distance=8mm and 12mm,
  block/.style={
    rectangle, rounded corners=2pt, draw=blue!60!black, fill=blue!8,
    minimum width=42mm, minimum height=9mm, align=center, font=\small
  },
  data/.style={
    rectangle, rounded corners=1pt, draw=gray!60, fill=gray!10,
    minimum width=42mm, minimum height=8mm, align=center, font=\footnotesize\ttfamily
  },
  llm/.style={
    rectangle, rounded corners=2pt, draw=red!60!black, fill=red!8,
    minimum width=42mm, minimum height=9mm, align=center, font=\small
  },
  store/.style={
    cylinder, shape border rotate=90, aspect=0.25,
    draw=teal!60!black, fill=teal!10, minimum width=30mm, minimum height=14mm,
    align=center, font=\small
  },
  doc/.style={
    rectangle, rounded corners=1pt, draw=orange!70!black, fill=orange!8,
    minimum width=36mm, minimum height=7mm, align=center, font=\scriptsize
  },
  parallel/.style={
    rectangle, rounded corners=2pt, draw=violet!60!black, fill=violet!8,
    minimum width=48mm, minimum height=9mm, align=center, font=\small
  },
  flow/.style={-{Stealth[length=2.2mm]}, thick, draw=blue!60!black},
  flowdata/.style={-{Stealth[length=2.2mm]}, thick, draw=gray!60, dashed},
  label/.style={font=\scriptsize\itshape, fill=white, inner sep=1pt},
}

% ---------- Titre ----------
\title{\textbf{Phase 2 — RAG sur la documentation Python}\\
\large Cinq concepts essentiels + six exercices pour aller plus loin}
\author{Cours PyTutor2}
\date{}

\begin{document}
\maketitle

\begin{abstract}
\noindent
Cette phase 2 construit, brique par brique, un système de
\emph{Retrieval-Augmented Generation} (RAG) qui répond à des questions
en s'appuyant sur la documentation Python officielle. On y aborde~:
l'indexation (Web $\to$ embeddings $\to$ Chroma persistante), le
\texttt{Retriever} LCEL, la chaîne RAG canonique en LCEL, la récupération
des sources avec \texttt{RunnableParallel}, et la comparaison
\emph{similarity vs MMR}. Chaque exercice suit la même structure que
la Phase 1 (énoncé + objectifs en boîte, solution détaillée, vérification
LangSmith, récapitulatif). Une seconde partie
\emph{« Pour aller plus loin »} prolonge avec six exercices RAG avancés~:
impact du \texttt{chunk\_size}, mode interactif CLI, prompt strict
anti-hallucination, observabilité avec LangSmith,
\texttt{MultiQueryRetriever}, et filtrage par métadonnée.
\end{abstract}

\tableofcontents
\bigskip

\part{Exercices}

% #####################################################################
% #####################################################################
\section{Indexation : du Web vers la base vectorielle}
% #####################################################################
% #####################################################################

\begin{tcolorbox}[title=Exercice 1. Construire un index Chroma persistant]
Aspire 7 pages de la documentation officielle Python
(\texttt{docs.python.org/3/tutorial/*}), découpe-les en \emph{chunks},
encode-les en vecteurs avec un modèle d'embeddings local, et persiste
l'index sur disque dans un dossier \texttt{chroma\_db/}.

Tu écriras une fonction \texttt{build\_or\_load\_vectorstore(force\_rebuild=False)}
qui~: (a) \textbf{recharge} l'index depuis le disque si \texttt{chroma\_db/}
existe ; (b) \textbf{reconstruit} l'index si l'argument \texttt{force\_rebuild}
est vrai ou si le dossier n'existe pas.

\textbf{Contrainte}~: l'indexation initiale prend $\sim$30 s (téléchargement
des embeddings + scraping des URL) ; les lancements suivants doivent
prendre $<$1 s grâce à la persistance.

\bigskip
\noindent\textbf{Objectifs pédagogiques}
\begin{itemize}[leftmargin=*]
  \item Utiliser \texttt{WebBaseLoader} pour aspirer une liste d'URL et
        obtenir des \texttt{Document} LangChain (texte + \texttt{metadata}).
  \item Choisir un \texttt{chunk\_size} et un \texttt{chunk\_overlap}
        cohérents pour de la doc technique, avec
        \texttt{RecursiveCharacterTextSplitter}.
  \item Encoder du texte en vecteurs avec
        \texttt{HuggingFaceEmbeddings} (modèle local, gratuit).
  \item Persister un \texttt{Chroma} sur disque et savoir le recharger
        sans réindexer.
\end{itemize}
\end{tcolorbox}

\subsection{Solution}

\subsubsection{Architecture}

\begin{center}
\begin{tikzpicture}[node distance=5mm and 14mm]
  \node[doc] (urls) {7 URL\\docs.python.org};
  \node[block, below=of urls] (loader) {WebBaseLoader\\\scriptsize \texttt{.load()}};
  \node[doc, below=of loader] (raw) {7 \texttt{Document}\\(une page = un doc)};
  \node[block, below=of raw] (splitter) {RecursiveCharacterTextSplitter\\\scriptsize \texttt{chunk\_size=1000, overlap=200}};
  \node[doc, below=of splitter] (chunks) {$\sim$300 chunks};
  \node[block, below=of chunks] (embeds) {HuggingFaceEmbeddings\\\scriptsize \texttt{all-MiniLM-L6-v2} (384-d)};
  \node[store, below=of embeds] (chroma) {Chroma\\\scriptsize \texttt{chroma\_db/}};

  \draw[flow] (urls)    -- (loader);
  \draw[flow] (loader)  -- (raw);
  \draw[flow] (raw)     -- (splitter);
  \draw[flow] (splitter)-- (chunks);
  \draw[flow] (chunks)  -- (embeds);
  \draw[flow] (embeds)  -- (chroma) node[label, midway, right]{persistance};
\end{tikzpicture}
\end{center}

Un pipeline en quatre étapes~: \textbf{scraping}, \textbf{découpage},
\textbf{encodage}, \textbf{stockage}. La séparation entre indexation
(lente, faite une fois) et interrogation (rapide, à chaque question)
est essentielle pour un RAG utilisable au quotidien.

\subsubsection{Le code complet}

\begin{pythoncode}
from pathlib import Path

from langchain_community.document_loaders import WebBaseLoader
from langchain_text_splitters import RecursiveCharacterTextSplitter
from langchain_huggingface import HuggingFaceEmbeddings
from langchain_chroma import Chroma


PYTHON_DOCS_URLS = [
    "https://docs.python.org/3/tutorial/controlflow.html",
    "https://docs.python.org/3/tutorial/datastructures.html",
    "https://docs.python.org/3/tutorial/classes.html",
    "https://docs.python.org/3/tutorial/errors.html",
    "https://docs.python.org/3/tutorial/modules.html",
    "https://docs.python.org/3/howto/functional.html",
    "https://docs.python.org/3/glossary.html",
]

PERSIST_DIR     = "./chroma_db"
EMBEDDING_MODEL = "sentence-transformers/all-MiniLM-L6-v2"
COLLECTION_NAME = "python_docs"


def get_embeddings() -> HuggingFaceEmbeddings:
    """Singleton de fait : un seul modèle d'embeddings réutilisé partout."""
    return HuggingFaceEmbeddings(model_name=EMBEDDING_MODEL)


def build_or_load_vectorstore(force_rebuild: bool = False) -> Chroma:
    """Recharge l'index si chroma_db/ existe, sinon le construit."""
    embeddings = get_embeddings()
    persist_path = Path(PERSIST_DIR)

    if persist_path.exists() and not force_rebuild:
        print(f"✓ Rechargement de l'index depuis {PERSIST_DIR}")
        return Chroma(
            collection_name=COLLECTION_NAME,
            persist_directory=PERSIST_DIR,
            embedding_function=embeddings,
        )

    # --- 1. Scraping : URL → Document ---
    loader = WebBaseLoader(PYTHON_DOCS_URLS)
    loader.requests_kwargs = {"headers": {"User-Agent": "PyTutor-Edu/0.1"}}
    raw_docs = loader.load()

    # --- 2. Découpage : Document → chunks ---
    splitter = RecursiveCharacterTextSplitter(
        chunk_size=1000,
        chunk_overlap=200,
        separators=["\n\n", "\n", ". ", " ", ""],
    )
    chunks = splitter.split_documents(raw_docs)

    # --- 3. Encodage + 4. Persistance (en une étape) ---
    vectorstore = Chroma.from_documents(
        documents=chunks,
        embedding=embeddings,
        collection_name=COLLECTION_NAME,
        persist_directory=PERSIST_DIR,
    )
    return vectorstore
\end{pythoncode}

\subsubsection{Explications pas à pas}

\subsubsection*{Étape 1 — \texttt{WebBaseLoader} : du Web au \texttt{Document}}

\texttt{WebBaseLoader(urls)} accepte une URL ou une liste d'URL. Chaque
page devient un \texttt{langchain\_core.documents.Document}~:

\begin{pythoncode}
Document(
    page_content="Texte complet de la page...",
    metadata={"source": "https://docs.python.org/3/tutorial/classes.html",
              "title": "9. Classes"},
)
\end{pythoncode}

Le champ \texttt{metadata["source"]} est précieux~: il permettra plus tard
de \textbf{citer les sources} et de filtrer.

\begin{quote}\itshape
\textbf{User-Agent personnalisé}~:
\texttt{loader.requests\_kwargs = \{"headers": \{"User-Agent": "PyTutor-Edu/0.1"\}\}}.
Certains sites bloquent les bots anonymes ; mieux vaut s'identifier
explicitement.
\end{quote}

\subsubsection*{Étape 2 — \texttt{RecursiveCharacterTextSplitter} : le découpage}

Un \texttt{Document} de 50~000 caractères est inutilisable directement~:
trop gros pour le contexte du LLM, et la recherche par similarité perd
en précision sur des blocs trop larges. On découpe en \emph{chunks}~:

\begin{pythoncode}
RecursiveCharacterTextSplitter(
    chunk_size=1000,                                # en CARACTÈRES
    chunk_overlap=200,                              # 20% de chevauchement
    separators=["\n\n", "\n", ". ", " ", ""],       # ordre prioritaire
)
\end{pythoncode}

Trois paramètres clés~:

\begin{itemize}[leftmargin=*]
  \item \textbf{\texttt{chunk\_size=1000}}~: en caractères, \emph{pas en
        tokens}. 1000 caractères $\approx$ 200--250 tokens — bon compromis
        pour de la doc technique.
  \item \textbf{\texttt{chunk\_overlap=200}}~: 20\% de chevauchement entre
        chunks consécutifs. Évite de couper une idée en deux.
  \item \textbf{\texttt{separators}}~: l'ordre compte. Le splitter essaie
        d'abord \texttt{"\textbackslash n\textbackslash n"} (paragraphe),
        puis \texttt{"\textbackslash n"} (ligne), puis \texttt{". "} (phrase),
        et en dernier recours au milieu d'un mot.
\end{itemize}

\begin{quote}\itshape
La \emph{taille du chunk} est l'hyperparamètre n\textsuperscript{o}1 de
qualité d'un RAG. Trop petit~: contexte fragmenté, le modèle manque le
sens global. Trop grand~: le retriever distingue mal les passages
pertinents des passages tangentiels. 800-1200 caractères est le sweet
spot pour de la doc technique.
\end{quote}

\subsubsection*{Étape 3 — \texttt{HuggingFaceEmbeddings} : encoder en vecteurs}

L'embedding transforme une chaîne de caractères en vecteur dense de
dimension fixe ($\sim$384 ici). Deux passages sémantiquement proches
donnent des vecteurs proches en distance cosinus.

\begin{pythoncode}
HuggingFaceEmbeddings(model_name="sentence-transformers/all-MiniLM-L6-v2")
\end{pythoncode}

Le modèle \texttt{all-MiniLM-L6-v2}~:
\begin{itemize}[leftmargin=*]
  \item \textbf{384 dimensions} (compact, rapide)
  \item \textbf{Multilingue} (utile pour des questions en français sur
        de la doc anglaise)
  \item \textbf{Local} (pas d'appel API, pas de coût par requête)
  \item \textbf{$\sim$80~MB} (téléchargé une fois, mis en cache)
\end{itemize}

\subsubsection*{Étape 4 — \texttt{Chroma} : stocker et rechercher}

\texttt{Chroma} est une base de données vectorielle qui~:
\begin{itemize}[leftmargin=*]
  \item Indexe les vecteurs pour la recherche par similarité.
  \item Stocke \texttt{page\_content} et \texttt{metadata} à côté.
  \item Persiste dans un dossier sur disque
        (\texttt{persist\_directory}).
\end{itemize}

\texttt{Chroma.from\_documents(\dots)} fait tout en une étape~: calcule
les embeddings, indexe, persiste. Pour recharger plus tard~:

\begin{pythoncode}
Chroma(
    collection_name="python_docs",
    persist_directory="./chroma_db",
    embedding_function=embeddings,  # même modèle qu'à l'indexation !
)
\end{pythoncode}

\subsubsection*{Étape 5 — Le pattern \emph{lazy rebuild}}

L'indexation est lente. On la fait une fois et on recharge ensuite~:

\begin{pythoncode}
if persist_path.exists() and not force_rebuild:
    return Chroma(...)             # rechargement <1s
# sinon, indexation complète
\end{pythoncode}

C'est le pattern à connaître pour tous les artefacts coûteux à construire~:
embeddings, modèles fine-tunés, caches de prompts, etc.

\begin{tcolorbox}[
  colback=white,
  colframe=black!50,
  coltitle=black,
  fonttitle=\footnotesize\bfseries,
  fontupper=\footnotesize,
  title=Sortie de l'exécution,
  breakable,
]
\begin{verbatim}
-> Indexation de 7 pages de docs.python.org...
   7 pages chargées
   327 chunks générés
[OK] Index créé et persisté dans ./chroma_db
\end{verbatim}
\end{tcolorbox}

\subsection{Vérification dans LangSmith}

L'indexation pure n'implique \textbf{pas} d'appel LLM, donc rien à voir
dans LangSmith ici. Tu peux vérifier le contenu de l'index directement~:

\begin{pythoncode}
vs = build_or_load_vectorstore()
print(vs._collection.count())          # nombre de chunks indexés
docs = vs.similarity_search("lambda", k=3)
for d in docs:
    print(d.metadata["source"], d.page_content[:100])
\end{pythoncode}

\begin{quote}\itshape
LangSmith ne trace que les \texttt{Runnable} (chaînes, appels LLM,
retrievers). L'indexation Chroma se passe \emph{en dehors} de ce
périmètre — c'est pourquoi on l'isole soigneusement du reste du code.
\end{quote}

\subsection*{Récapitulatif}
{\itshape\sloppy
\textit{Ce qu'il faut retenir :}
\begin{enumerate}[leftmargin=*]
  \item Le pipeline d'indexation est universel~: \texttt{Loader $\to$
        Splitter $\to$ Embeddings $\to$ VectorStore}.
  \item \texttt{chunk\_size} est le \textbf{premier} hyperparamètre à
        régler. 800--1200 caractères pour de la doc technique.
  \item \texttt{chunk\_overlap} $\approx$ 20\% du \texttt{chunk\_size}
        évite de couper une idée en deux.
  \item \texttt{HuggingFaceEmbeddings} est local, gratuit, multilingue.
        L'alternative \texttt{OpenAIEmbeddings} est plus performante mais
        payante.
  \item \texttt{Chroma} persiste sur disque. Pattern \emph{lazy rebuild}~:
        recharger si l'index existe, sinon reconstruire.
\end{enumerate}
\par}

% #####################################################################
% #####################################################################
\section{Retriever : recherche par similarité}
% #####################################################################
% #####################################################################

\begin{tcolorbox}[title=Exercice 2. Transformer un vectorstore en \texttt{Runnable} LCEL]
À partir d'un \texttt{Chroma} indexé, construis un \texttt{Retriever}
qui prend une chaîne de caractères en entrée et renvoie les \texttt{k}
documents les plus pertinents par similarité cosinus.

\textbf{Contrainte}~: le \texttt{Retriever} doit être un
\texttt{Runnable} LCEL, c'est-à-dire qu'il doit pouvoir être composé
avec \texttt{|}, batché, streamé, et tracé dans LangSmith comme n'importe
quel autre maillon.

\bigskip
\noindent\textbf{Objectifs pédagogiques}
\begin{itemize}[leftmargin=*]
  \item Comprendre comment \texttt{vectorstore.as\_retriever()} produit
        un \texttt{Runnable} à partir d'un index Chroma.
  \item Paramétrer la recherche~: \texttt{k} (nombre de docs retournés),
        \texttt{search\_type} (\texttt{"similarity"} par défaut).
  \item Savoir lire ce qui se passe sous le capot~: \emph{embedding de
        la query $\to$ produit scalaire $\to$ tri $\to$ top-k}.
\end{itemize}
\end{tcolorbox}

\subsection{Solution}

\subsubsection{Architecture}

\begin{center}
\resizebox{\linewidth}{!}{%
\begin{tikzpicture}[node distance=6mm and 14mm]
  \node[data] (query) {"Comment écrire un\\générateur en Python ?"};
  \node[block, right=of query] (embed) {Embeddings\\\scriptsize encode la query};
  \node[data, right=of embed] (qvec) {vecteur (384-d)};

  \node[store, below=12mm of embed] (chroma) {Chroma\\$\sim$300 vecteurs};
  \node[block, below=12mm of qvec] (sim) {Recherche\\\scriptsize cosine + top-k};

  \node[doc, right=of sim] (docs) {[Doc, Doc, Doc, Doc]\\\scriptsize (k=4)};

  \draw[flowdata] (query) -- (embed);
  \draw[flow]     (embed) -- (qvec);
  \draw[flow]     (qvec)  -- (sim);
  \draw[flow]     (chroma)-- (sim) node[label, midway, above, sloped]{candidats};
  \draw[flow]     (sim)   -- (docs);
\end{tikzpicture}%
}
\end{center}

Le retriever fait deux choses~: (1) encoder la \emph{query} avec le même
modèle d'embeddings que les documents, (2) chercher les \texttt{k}
vecteurs les plus proches par similarité cosinus.

\subsubsection{Le code complet}

\begin{pythoncode}
# Le retriever est obtenu d'une seule ligne à partir du vectorstore.
vectorstore = build_or_load_vectorstore()

retriever = vectorstore.as_retriever(
    search_type="similarity",        # défaut, on l'écrit pour la clarté
    search_kwargs={"k": 4},          # nombre de docs renvoyés
)

# Le retriever est un Runnable : .invoke, .batch, .stream marchent
docs = retriever.invoke("Comment écrire un générateur en Python ?")

for i, d in enumerate(docs, 1):
    print(f"[{i}] {d.metadata['source']}")
    print(d.page_content[:200], "...")
\end{pythoncode}

\subsubsection{Explications pas à pas}

\subsubsection*{Étape 1 — \texttt{.as\_retriever()} : du store au \texttt{Runnable}}

Un \texttt{Chroma} (ou n'importe quel \texttt{VectorStore} LangChain)
n'est pas directement un \texttt{Runnable}. La méthode
\texttt{.as\_retriever()} en crée un, qui peut alors être composé dans
une chaîne LCEL avec l'opérateur \texttt{|}.

\begin{pythoncode}
retriever = vectorstore.as_retriever(...)
# retriever est un Runnable
type(retriever).__mro__   # contient Runnable[str, list[Document]]
\end{pythoncode}

Le contrat de type~: \texttt{str $\to$ list[Document]}.

\subsubsection*{Étape 2 — Ce qui se passe à l'\texttt{.invoke()}}

\begin{enumerate}[leftmargin=*]
  \item La \texttt{query} (str) est encodée avec le \textbf{même
        \texttt{HuggingFaceEmbeddings}} qui a indexé les documents.
        C'est crucial~: utiliser un autre modèle d'embeddings produirait
        un vecteur incomparable avec ceux de l'index.
  \item Chroma calcule la similarité cosinus entre le vecteur query et
        \textbf{tous} les vecteurs indexés (en pratique avec un index
        approximatif pour aller vite).
  \item Les \texttt{k} chunks de plus haute similarité sont renvoyés,
        triés par score décroissant.
\end{enumerate}

\begin{quote}\itshape
\textbf{Pourquoi cosinus et pas euclidienne~?} La distance euclidienne
dépend de la norme des vecteurs ; la similarité cosinus ne dépend que de
l'angle. Pour des embeddings de phrases, c'est l'angle qui porte le sens.
\end{quote}

\subsubsection*{Étape 3 — Le paramètre \texttt{k}}

\texttt{k} est le nombre de chunks à renvoyer. C'est l'hyperparamètre
n\textsuperscript{o}2 de qualité du RAG (après \texttt{chunk\_size})~:

\begin{itemize}[leftmargin=*]
  \item \texttt{k=1}~: trop pauvre, on rate souvent du contexte adjacent.
  \item \texttt{k=4} à \texttt{k=6}~: standard pour de la doc technique.
  \item \texttt{k=10+}~: contexte saturé, le LLM peut perdre le fil.
        Coût en tokens augmenté.
\end{itemize}

\subsubsection*{Étape 4 — Le retriever dans une chaîne LCEL}

Comme tout \texttt{Runnable}, le retriever se compose~:

\begin{pythoncode}
# Construire un mini-pipeline
def truncate(docs, n=3):
    return docs[:n]

mini_chain = retriever | truncate
# Mini-chain reçoit une str et renvoie les 3 premiers docs
\end{pythoncode}

On verra dans l'exercice 3 comment l'enchaîner avec un prompt et un
modèle pour faire un RAG complet.

\begin{tcolorbox}[
  colback=white,
  colframe=black!50,
  coltitle=black,
  fonttitle=\footnotesize\bfseries,
  fontupper=\footnotesize,
  title=Sortie de l'exécution,
  breakable,
]
\begin{verbatim}
[1] https://docs.python.org/3/tutorial/modules.html
'UnicodeError', 'UnicodeTranslateError', 'UnicodeWarning', 'UserWarning',
'ValueError', 'Warning', 'ZeroDivisionError', '_', '__build_class__',
'__debug__', '__doc__', '__import__', '__name__', '__pac ...
[2] https://docs.python.org/3/howto/functional.html
GeneratorExit
occurs, I suggest
using a
try:
...
finally:
suite instead of catching
GeneratorExit
.
The cumulative effect of these changes is to turn generators from one-way
producers of information i ...
[3] https://docs.python.org/3/tutorial/classes.html
9. Classes
9.1. A Word About Names and Objects
9.2. Python Scopes and Namespaces
9.2.1. Scopes and Namespaces Example
9.3. A First Look at Classes
9.3.1. Class Definition Syntax
9.3.2. Class Objects
9 ...
[4] https://docs.python.org/3/howto/functional.html
this case. Don't just use its value in expressions unless you're sure that
the send() method will be the only method used to resume your generator
function. In addition to send(), there are two othe ...
\end{verbatim}
\end{tcolorbox}

\subsection{Vérification dans LangSmith}

Quand le retriever est utilisé dans une chaîne tracée, LangSmith affiche~:

\begin{center}
\begin{tikzpicture}[every node/.style={font=\footnotesize\ttfamily, anchor=west}]
  \node                              (a) {Retriever (VectorStoreRetriever)};
  \node[below=1mm of a, xshift=4mm]  (b) {Input  : "Comment écrire un générateur en Python ?"};
  \node[below=1mm of b]              (c) {Output : [Document, Document, Document, Document]};
  \node[below=1mm of c, xshift=4mm]  (d) {[0] source: tutorial/classes.html};
  \node[below=1mm of d]              (e) {[1] source: howto/functional.html};
  \node[below=1mm of e]              (f) {[2] ...};
\end{tikzpicture}
\end{center}

\begin{quote}\itshape
L'onglet \emph{Output} montre la liste des \texttt{Document} renvoyés
avec leur \texttt{metadata} (notamment \texttt{source}). C'est le
\textbf{premier endroit où regarder} quand un RAG répond mal~: souvent
le retriever a ramené du contexte hors-sujet.
\end{quote}

\subsection*{Récapitulatif}
{\itshape\sloppy
\textit{Ce qu'il faut retenir :}
\begin{enumerate}[leftmargin=*]
  \item \texttt{vectorstore.as\_retriever(\dots)} transforme un store
        en \texttt{Runnable} LCEL, composable avec \texttt{|}.
  \item Contrat du retriever~: \texttt{str $\to$ list[Document]}.
  \item Sous le capot~: encode la query avec le \emph{même} modèle
        d'embeddings, cherche les top-\texttt{k} par cosinus.
  \item \texttt{k=4} à \texttt{k=6} est le standard pour de la doc
        technique.
  \item Premier point de débogage d'un RAG cassé~: regarder l'output du
        retriever dans LangSmith.
\end{enumerate}
\par}

% #####################################################################
% #####################################################################
\section{Chaîne RAG simple en LCEL}
% #####################################################################
% #####################################################################

\begin{tcolorbox}[title=Exercice 3. Composer un RAG end-to-end en LCEL]
Construis une chaîne \texttt{Runnable} qui prend une \textbf{question
en texte brut} et renvoie une \textbf{réponse en texte brut}, en
s'appuyant sur le retriever de l'exercice 2 pour récupérer du contexte
pertinent dans la doc Python.

\textbf{Contrainte}~: utiliser le pattern canonique LCEL avec un dict
implicite qui contient deux branches \texttt{context} et \texttt{question},
puis composer \texttt{| prompt | model | StrOutputParser()}. Le modèle
doit \textbf{citer ses sources} par numéro \texttt{[N]}.

\bigskip
\noindent\textbf{Objectifs pédagogiques}
\begin{itemize}[leftmargin=*]
  \item Reconnaître le \textbf{pattern canonique RAG en LCEL}~:
        \texttt{\{context, question\} | prompt | model | parser}.
  \item Utiliser \texttt{RunnablePassthrough()} pour faire passer la
        question à travers une branche du dict implicite.
  \item Composer \texttt{retriever | format\_docs} pour transformer la
        liste de \texttt{Document} en string utilisable dans le prompt.
  \item Donner au modèle des instructions claires pour s'appuyer sur
        le contexte fourni et citer les sources.
\end{itemize}
\end{tcolorbox}

\subsection{Solution}

\subsubsection{Architecture}

\begin{center}
\begin{tikzpicture}[node distance=5mm and 14mm]
  \node[data] (q) {"Comment écrire un\\générateur en Python ?"};

  \node[parallel, below=10mm of q] (dict) {\{context, question\}\\\scriptsize (RunnableParallel implicite)};

  \node[block, below left=10mm and 12mm of dict]  (ret)  {retriever\\\scriptsize | format\_docs};
  \node[block, below right=10mm and 12mm of dict] (pass) {RunnablePassthrough\\\scriptsize (question inchangée)};

  \node[block, below=22mm of dict] (prompt) {RAG\_PROMPT};
  \node[llm,   below=of prompt]    (model)  {Modèle de chat};
  \node[block, below=of model]     (parser) {StrOutputParser};
  \node[data,  below=of parser]    (out)    {réponse (str)};

  \draw[flowdata] (q)    -- (dict);
  \draw[flow]     (dict) -- (ret);
  \draw[flow]     (dict) -- (pass);
  \draw[flow]     (ret)  -- (prompt) node[label, midway, sloped, above]{context: "[1] ..."};
  \draw[flow]     (pass) -- (prompt) node[label, midway, sloped, above]{question};
  \draw[flow]     (prompt)-- (model);
  \draw[flow]     (model) -- (parser);
  \draw[flow]     (parser)-- (out);
\end{tikzpicture}
\end{center}

L'input (str) est dupliqué dans deux branches qui s'exécutent en
parallèle~: l'une va chercher le contexte, l'autre laisse passer la
question. Les deux résultats sont assemblés dans un dict que le prompt
sait formater.

\subsubsection{Le code complet}

\begin{pythoncode}
from langchain_core.documents import Document
from langchain_core.prompts import ChatPromptTemplate
from langchain_core.output_parsers import StrOutputParser
from langchain_core.runnables import RunnablePassthrough
from config import get_model


RAG_PROMPT = ChatPromptTemplate.from_messages([
    (
        "system",
        "Tu es un assistant pédagogique pour des élèves de Python. "
        "Réponds à la question en t'appuyant UNIQUEMENT sur le contexte fourni.\n"
        "Si la réponse n'est pas dans le contexte, dis-le explicitement plutôt "
        "que d'inventer. Cite le ou les chunks utilisés par leur numéro entre "
        "crochets, par exemple [1] ou [2,3]."
    ),
    (
        "human",
        "Contexte :\n{context}\n\n"
        "Question : {question}\n\n"
        "Réponse pédagogique :"
    ),
])


def format_docs(docs: list[Document]) -> str:
    """Concatène les documents en un bloc de contexte numéroté."""
    return "\n\n".join(
        f"[{i + 1}] (source: {doc.metadata.get('source', '?')})\n{doc.page_content}"
        for i, doc in enumerate(docs)
    )


def build_rag_chain(retriever):
    """Chaîne RAG simple : str -> str."""
    return (
        {
            "context":  retriever | format_docs,
            "question": RunnablePassthrough(),
        }
        | RAG_PROMPT
        | get_model()
        | StrOutputParser()
    )


# Usage
vectorstore = build_or_load_vectorstore()
retriever   = vectorstore.as_retriever(search_kwargs={"k": 4})
chain       = build_rag_chain(retriever)

answer = chain.invoke("Comment écrire un générateur en Python ?")
print(answer)
\end{pythoncode}

\subsubsection{Explications pas à pas}

\subsubsection*{Étape 1 — Le dict implicite = \texttt{RunnableParallel}}

\begin{pythoncode}
{
    "context":  retriever | format_docs,
    "question": RunnablePassthrough(),
}
\end{pythoncode}

Quand un dict apparaît dans une chaîne LCEL (au début ou via
\texttt{|}), LCEL le convertit \textbf{implicitement} en
\texttt{RunnableParallel}. Les deux branches reçoivent le même input
(la question, str) et s'exécutent en parallèle~:

\begin{itemize}[leftmargin=*]
  \item \textbf{branche \texttt{context}}~: la string passe par
        \texttt{retriever} (qui renvoie une \texttt{list[Document]}),
        puis par \texttt{format\_docs} (qui renvoie une str).
  \item \textbf{branche \texttt{question}}~: la string traverse
        \texttt{RunnablePassthrough()} inchangée.
\end{itemize}

La sortie est un dict \texttt{\{"context": "[1] ...", "question": "..."\}}
prêt à être consommé par le prompt.

\begin{quote}\itshape
Pattern à mémoriser~: \textbf{\texttt{\{ "key1": branch1, "key2":
branch2 \} | prompt | model | parser}}. C'est la forme canonique du RAG
en LCEL. Tu la reverras 1000 fois.
\end{quote}

\subsubsection*{Étape 2 — \texttt{retriever | format\_docs}}

\texttt{retriever} produit une \texttt{list[Document]}. Le prompt attend
une chaîne de caractères pour la variable \texttt{\{context\}}. La
composition \texttt{retriever | format\_docs} chaîne les deux~:

\begin{pythoncode}
def format_docs(docs):
    return "\n\n".join(
        f"[{i+1}] (source: {d.metadata['source']})\n{d.page_content}"
        for i, d in enumerate(docs)
    )
\end{pythoncode}

Trois détails~:
\begin{itemize}[leftmargin=*]
  \item Numérotation \texttt{[1]}, \texttt{[2]}, \dots~: indispensable
        pour que le modèle puisse citer les sources \emph{par numéro}.
  \item Ajout de la \texttt{source} dans chaque bloc~: traçabilité, et
        permet au modèle de citer l'URL.
  \item Séparateur \texttt{"\textbackslash n\textbackslash n"}~: les LLM
        repèrent mieux les blocs distincts ainsi qu'avec un simple
        \texttt{"\textbackslash n"}.
\end{itemize}

\subsubsection*{Étape 3 — \texttt{RunnablePassthrough}}

\texttt{RunnablePassthrough()} est l'identité de LCEL~: il prend un input
et le renvoie inchangé. Utile dans un \texttt{RunnableParallel} pour
\emph{conserver} une valeur d'entrée tout en faisant autre chose avec
elle dans une autre branche.

\begin{pythoncode}
RunnablePassthrough().invoke("hello")   # → "hello"
\end{pythoncode}

\subsubsection*{Étape 4 — Le prompt RAG}

Le prompt système contient \textbf{trois consignes} essentielles~:

\begin{enumerate}[leftmargin=*]
  \item « UNIQUEMENT sur le contexte fourni » — anti-hallucination.
  \item « Si la réponse n'est pas dans le contexte, dis-le » — préserve
        l'honnêteté épistémique.
  \item « Cite les chunks par numéro \texttt{[N]} » — traçabilité.
\end{enumerate}

\begin{quote}\itshape
Les trois consignes du prompt RAG sont \emph{défensives}. Les LLM
modernes ont tendance à compléter avec leurs connaissances générales —
en RAG, c'est exactement ce qu'on veut éviter (sinon, à quoi sert le
retrieval~?). On verra dans l'exercice bonus 3 comment durcir encore
le prompt.
\end{quote}

\subsubsection*{Étape 5 — \texttt{StrOutputParser} : extraire le texte}

Le modèle renvoie un \texttt{AIMessage}. Pour récupérer juste le texte
de la réponse, on ajoute \texttt{| StrOutputParser()} en fin de chaîne.
C'est le même parseur qu'on utilisait en streaming dans la Phase 1.

\begin{tcolorbox}[
  colback=white,
  colframe=black!50,
  coltitle=black,
  fonttitle=\footnotesize\bfseries,
  fontupper=\footnotesize,
  title=Sortie de l'exécution,
  breakable,
]
\begin{verbatim}
======================================================================
DEMO 1 -- RAG simple (retriever similarity, k=4)
======================================================================

Question : Comment écrire un générateur en Python ?

# Comment écrire un générateur en Python ?

D'après le contexte fourni, voici ce que je peux te dire sur les
générateurs :

## Définition
Un **générateur** est une fonction qui contient une ou plusieurs
expressions `yield` pour produire une série de valeurs [4].

## Caractéristiques principales
- Un générateur ressemble à une fonction normale, sauf qu'il utilise
  le mot-clé `yield` [4]
- Les valeurs produites peuvent être utilisées dans une boucle `for`
  ou récupérées une à une avec la fonction `next()` [4]

## Exemple de fonctionnement
Selon le contexte, quand tu appelles `next()` sur un générateur, il
exécute le code jusqu'à la prochaine expression `yield`. Quand il n'y
a plus de valeurs à produire, une exception `StopIteration` est
levée [1].

## Points importants à retenir
- Attention : `yield` retourne souvent `None`, donc tu dois vérifier
  cette valeur [1]
- Les générateurs ont des méthodes spéciales comme `throw()` et
  `close()` pour les contrôler [1]

**Cependant**, le contexte fourni ne contient pas d'exemple complet
de code montrant comment écrire un générateur pas à pas. Pour avoir
un exemple détaillé avec la syntaxe exacte, tu devrais consulter la
documentation Python officielle sur les générateurs.
\end{verbatim}
\end{tcolorbox}

\subsection{Vérification dans LangSmith}

L'arbre d'exécution attendu~:

\begin{center}
\begin{tikzpicture}[every node/.style={font=\footnotesize\ttfamily, anchor=west}]
  \node                              (a) {RunnableSequence};
  \node[below=1mm of a, xshift=4mm]  (b) {RunnableParallel<context, question>};
  \node[below=1mm of b, xshift=4mm]  (c) {[context] RunnableSequence};
  \node[below=1mm of c, xshift=4mm]  (d) {Retriever ← 1 query embedée};
  \node[below=1mm of d]              (e) {format\_docs (RunnableLambda)};
  \node[below=1mm of e, xshift=-4mm] (f) {[question] RunnablePassthrough};
  \node[below=1mm of f, xshift=-4mm] (g) {ChatPromptTemplate};
  \node[below=1mm of g]              (h) {ChatAnthropic ← appel LLM};
  \node[below=1mm of h]              (i) {StrOutputParser};
\end{tikzpicture}
\end{center}

Les trois points à regarder~:
\begin{enumerate}[leftmargin=*]
  \item \textbf{Output du \texttt{Retriever}}~: les chunks ramenés
        sont-ils pertinents~? Si non, problème côté indexation ou
        \texttt{k}.
  \item \textbf{Input du \texttt{ChatAnthropic}}~: le prompt rendu, avec
        le contexte complet. C'est la \emph{vérité terrain} de ce que le
        modèle a vu.
  \item \textbf{Output du \texttt{ChatAnthropic}}~: la réponse cite-t-elle
        bien les \texttt{[N]}~? Si non, le modèle a peut-être ignoré le
        contexte (à durcir avec le prompt strict de l'exercice bonus 3).
\end{enumerate}

\begin{quote}\itshape
Si la réponse est incorrecte, \textbf{ne blâme pas le modèle d'abord} :
regarde l'output du retriever. 80\% des bugs RAG viennent d'un mauvais
retrieval (chunk\_size mal réglé, query mal formulée, k trop petit).
\end{quote}

\subsection*{Récapitulatif}
{\itshape\sloppy
\textit{Ce qu'il faut retenir :}
\begin{enumerate}[leftmargin=*]
  \item Pattern canonique RAG~: \texttt{\{context, question\} | prompt |
        model | parser}.
  \item Le \textbf{dict implicite} devient un \texttt{RunnableParallel}
        à deux branches~: une qui va chercher (retriever | format), une
        qui passe la question.
  \item \texttt{format\_docs} numérote les chunks et expose les sources —
        condition pour que le modèle puisse citer.
  \item Le prompt système doit explicitement~: (1) restreindre au
        contexte, (2) autoriser le « je ne sais pas », (3) exiger les
        citations \texttt{[N]}.
  \item Pour debugger~: lire d'abord l'output du retriever dans
        LangSmith, ensuite le prompt rendu.
\end{enumerate}
\par}

% #####################################################################
% #####################################################################
\section{RAG avec sources : récupérer réponse \emph{et} contexte}
% #####################################################################
% #####################################################################

\begin{tcolorbox}[title=Exercice 4. Préserver les documents sources dans la sortie]
Modifie la chaîne RAG de l'exercice 3 pour qu'elle renvoie un
\textbf{dict} \texttt{\{"answer": str, "context": list[Document],
"question": str\}} au lieu de juste la réponse en string. Cela permet
d'afficher les sources à l'utilisateur et de \emph{vérifier} sur quels
chunks la réponse s'est appuyée.

\textbf{Contrainte}~: une seule invocation du retriever (pas deux). On
utilise \texttt{RunnablePassthrough.assign} pour \emph{enrichir} le dict
au fil de la chaîne sans perdre les clés précédentes.

\bigskip
\noindent\textbf{Objectifs pédagogiques}
\begin{itemize}[leftmargin=*]
  \item Distinguer \texttt{RunnableParallel} (remplace l'input par un
        dict) de \texttt{RunnablePassthrough.assign} (enrichit le dict).
  \item Construire une chaîne en \textbf{deux phases}~: récupération du
        contexte (1 appel retriever), puis génération de la réponse en
        gardant le contexte en mémoire.
  \item Comprendre le pattern \emph{fan-out / accumulate} de LCEL.
\end{itemize}
\end{tcolorbox}

\subsection{Solution}

\subsubsection{Architecture}

\begin{center}
\begin{tikzpicture}[node distance=5mm and 14mm]
  \node[data] (q) {"À quoi sert yield ?"};

  \node[parallel, below=of q] (retrieve)
       {RunnableParallel\\\scriptsize \{context=retriever, question=passthrough\}};

  \node[data, below=of retrieve] (mid)
       {\{context: [Doc, Doc, Doc],\\\ question: "À quoi sert yield ?"\}};

  \node[block, below=of mid] (assign)
       {RunnablePassthrough.assign(answer=\dots)};

  \node[data, below=of assign] (out)
       {\{context: [Doc, Doc, Doc],\\\ question: "À quoi sert yield ?",\\\ answer: "Le mot-clé yield..."\}};

  \draw[flowdata] (q)        -- (retrieve);
  \draw[flow]     (retrieve) -- (mid)    node[label, midway, right]{1 appel retriever};
  \draw[flowdata] (mid)      -- (assign);
  \draw[flow]     (assign)   -- (out)    node[label, midway, right]{1 appel LLM (ajoute "answer")};
\end{tikzpicture}
\end{center}

Deux étapes propres~: d'abord on récupère le contexte une seule fois
(branche \texttt{retrieve}), ensuite on ajoute la réponse au dict
existant (\texttt{.assign}).

\subsubsection{Le code complet}

\begin{pythoncode}
from langchain_core.runnables import RunnableParallel, RunnablePassthrough


def build_rag_chain_with_sources(retriever):
    """Chaîne RAG qui renvoie un dict {answer, context, question}."""

    # Étape 1 : récupère le contexte UNE FOIS et le passe en aval
    retrieve = RunnableParallel(
        context=retriever,
        question=RunnablePassthrough(),
    )

    # Étape 2 : à partir de {context: [Doc...], question: str},
    # construit la réponse en formatant les docs au passage.
    answer_chain = (
        RunnablePassthrough.assign(
            context=lambda x: format_docs(x["context"]),
        )
        | RAG_PROMPT
        | get_model()
        | StrOutputParser()
    )

    # .assign(answer=...) ajoute la clé "answer" sans toucher aux autres.
    return retrieve | RunnablePassthrough.assign(answer=answer_chain)


# Usage
chain  = build_rag_chain_with_sources(retriever)
result = chain.invoke("À quoi sert le mot-clé yield exactement ?")

print("RÉPONSE :")
print(result["answer"])

print("\nSOURCES :")
for i, doc in enumerate(result["context"], 1):
    print(f"  [{i}] {doc.metadata['source']}")
    print(f"      {doc.page_content[:120]}...")
\end{pythoncode}

\subsubsection{Explications pas à pas}

\subsubsection*{Étape 1 — Récupérer une seule fois}

\begin{pythoncode}
retrieve = RunnableParallel(
    context=retriever,
    question=RunnablePassthrough(),
)
\end{pythoncode}

À partir d'une question (str), ce \texttt{RunnableParallel} produit~:
\begin{itemize}[leftmargin=*]
  \item \texttt{context}~: la \texttt{list[Document]} renvoyée par le
        retriever (\textbf{1 appel}).
  \item \texttt{question}~: la string d'origine (passthrough).
\end{itemize}

\textbf{Important}~: on ne formate \emph{pas encore} les docs ici. On
veut les garder bruts pour pouvoir les retourner à l'utilisateur en
sortie.

\subsubsection*{Étape 2 — La sous-chaîne qui calcule la réponse}

\begin{pythoncode}
answer_chain = (
    RunnablePassthrough.assign(
        context=lambda x: format_docs(x["context"]),
    )
    | RAG_PROMPT | get_model() | StrOutputParser()
)
\end{pythoncode}

Cette sous-chaîne reçoit \texttt{\{context: [Doc\dots], question: str\}}
et fait trois choses~:

\begin{enumerate}[leftmargin=*]
  \item \textbf{Remplace} la clé \texttt{context} par sa version formatée
        en string. Le \texttt{RunnablePassthrough.assign(context=\dots)}
        applique la lambda à l'input courant et place le résultat dans
        \texttt{context}. La clé \texttt{question} est intacte.
  \item Le prompt formate avec les variables \texttt{\{context\}} (str
        maintenant) et \texttt{\{question\}}.
  \item Modèle + parser produisent la réponse en string.
\end{enumerate}

\begin{quote}\itshape
\texttt{RunnablePassthrough.assign(key=fn)} \textbf{remplace} la clé
\texttt{key} si elle existe déjà, ou la \textbf{crée} si elle n'existe
pas. C'est l'équivalent d'un \texttt{UPDATE} en SQL sur un dict.
\end{quote}

\subsubsection*{Étape 3 — L'accumulation finale}

\begin{pythoncode}
return retrieve | RunnablePassthrough.assign(answer=answer_chain)
\end{pythoncode}

Le \texttt{RunnablePassthrough.assign(answer=answer\_chain)} prend le
dict produit par \texttt{retrieve} (\texttt{\{context: [Doc\dots],
question: str\}}) et y ajoute la clé \texttt{answer} (résultat de
\texttt{answer\_chain}).

La sortie finale~:

\begin{pythoncode}
{
    "context":  [Document, Document, Document, Document],  # bruts, intacts
    "question": "À quoi sert yield ?",                      # str
    "answer":   "Le mot-clé yield...",                      # str
}
\end{pythoncode}

L'utilisateur peut afficher \texttt{result["answer"]} et énumérer
\texttt{result["context"]} pour montrer les sources.

\subsubsection*{Étape 4 — Pourquoi pas deux retrievers~?}

Une approche naïve appellerait deux fois le retriever~: une fois pour
afficher les sources, une fois inclus dans la chaîne RAG. C'est deux
fois plus de calcul et l'index pourrait même retourner des docs
différents (selon la non-déterminisme).

Le pattern \texttt{RunnableParallel} + \texttt{.assign}~:
\begin{itemize}[leftmargin=*]
  \item \textbf{1 seul appel} retriever (économie).
  \item \textbf{Cohérence}~: la réponse s'appuie \emph{exactement} sur
        les sources affichées (pas de divergence possible).
  \item \textbf{Composable}~: tout reste un \texttt{Runnable}.
\end{itemize}

\begin{tcolorbox}[
  colback=white,
  colframe=black!50,
  coltitle=black,
  fonttitle=\footnotesize\bfseries,
  fontupper=\footnotesize,
  title=Sortie de l'exécution,
  breakable,
]
\begin{verbatim}
======================================================================
DEMO 2 -- RAG avec inspection des sources
======================================================================

Question : A quoi sert le mot-clé `yield` exactement ?

REPONSE :
# A quoi sert le mot-clé `yield` ?

Le mot-clé `yield` sert à **créer des générateurs** en Python. Voici
ses principales fonctions [1,2,3] :

## 1. **Suspendre l'exécution et retourner une valeur**
Quand `yield` est rencontré, la fonction s'arrête et retourne une
valeur. A chaque appel suivant, l'exécution reprend à partir du
`yield` [2].

## 2. **Préserver l'état local**
Contrairement à `return` qui termine la fonction, `yield` **suspend
l'exécution et préserve les variables locales**. Cela permet à la
fonction de reprendre son travail plus tard [2].

## 3. **Recevoir des valeurs**
`yield` peut aussi **recevoir des valeurs** via la méthode
`send(value)`. La valeur envoyée devient le résultat de
l'expression `yield` [1,3].

## Exemple simple [3] :
```python
def counter(maximum):
    i = 0
    while i < maximum:
        val = (yield i)  # Retourne i, puis attend une valeur
        if val is not None:
            i = val      # Recoit une nouvelle valeur
        else:
            i += 1
```

**En résumé** : `yield` transforme une fonction ordinaire en
**générateur** qui produit une séquence de valeurs une à une, tout en
gardant son état entre chaque appel. C'est très utile pour économiser
la mémoire avec de grandes séquences [2].

SOURCES UTILISEES :
  [1] https://docs.python.org/3/howto/functional.html
      are messy. In Python 2.5 there's a simple way to pass values into
      a generator. yield became an expression, returning a v...
  [2] https://docs.python.org/3/howto/functional.html
      yield and a return statement is that on reaching a yield the
      generator's state of execution is suspended and local varia...
  [3] https://docs.python.org/3/howto/functional.html
      + 12 .) Values are sent into a generator by calling its
      send(value) method.  This method resumes the generator's code an...
  [4] https://docs.python.org/3/howto/functional.html
      GeneratorExit occurs, I suggest using a try: ... finally: suite
      instead of catching GeneratorExit . The cumulative effec...
\end{verbatim}
\end{tcolorbox}

\subsection{Vérification dans LangSmith}

L'arbre d'exécution est nettement plus riche~:

\begin{center}
\begin{tikzpicture}[every node/.style={font=\footnotesize\ttfamily, anchor=west}]
  \node                              (a) {RunnableSequence};
  \node[below=1mm of a, xshift=4mm]  (b) {RunnableParallel<context, question>};
  \node[below=1mm of b, xshift=4mm]  (c) {[context] Retriever};
  \node[below=1mm of c]              (d) {[question] RunnablePassthrough};
  \node[below=1mm of d, xshift=-4mm] (e) {RunnableAssign<answer>};
  \node[below=1mm of e, xshift=4mm]  (f) {RunnableAssign<context=format\_docs>};
  \node[below=1mm of f]              (g) {ChatPromptTemplate};
  \node[below=1mm of g]              (h) {ChatAnthropic ← appel LLM};
  \node[below=1mm of h]              (i) {StrOutputParser};
\end{tikzpicture}
\end{center}

\begin{quote}\itshape
Dans l'onglet \emph{Output} du run racine, tu vois directement le dict
final avec ses 3 clés. C'est la matérialisation du pattern \emph{fan-out
/ accumulate} qu'on cherche en LCEL pour produire des sorties riches
sans dupliquer les calculs.
\end{quote}

\subsection*{Récapitulatif}
{\itshape\sloppy
\textit{Ce qu'il faut retenir :}
\begin{enumerate}[leftmargin=*]
  \item \texttt{RunnableParallel} crée un dict ; \texttt{Passthrough.assign}
        \textbf{enrichit} un dict existant. Ne pas confondre.
  \item Pattern \emph{fan-out / accumulate}~: une étape construit un
        dict, les étapes suivantes ajoutent des clés sans détruire les
        précédentes.
  \item Une seule invocation du retriever — la réponse et les sources
        affichées sont cohérentes par construction.
  \item Pour faire un RAG « professionnel » avec sources, c'est
        \textbf{le} pattern à connaître.
\end{enumerate}
\par}

% #####################################################################
% #####################################################################
\section{Recherche \emph{similarity} vs \emph{MMR}}
% #####################################################################
% #####################################################################

\begin{tcolorbox}[title=Exercice 5. Comparer deux stratégies de retrieval]
Compare deux \texttt{Retriever} sur la même question~: l'un avec
\texttt{search\_type="similarity"} (par défaut), l'autre avec
\texttt{search\_type="mmr"} (\emph{Maximal Marginal Relevance}).
Observe que MMR produit des chunks plus \emph{diversifiés} pour une
question généraliste comme \emph{« liste vs tuple »}.

\textbf{Contrainte}~: utiliser le même \texttt{Chroma}, juste changer
\texttt{search\_kwargs}. Pour MMR, expérimenter avec \texttt{fetch\_k}
(taille de la pool de candidats) et \texttt{lambda\_mult} (équilibre
similarité/diversité).

\bigskip
\noindent\textbf{Objectifs pédagogiques}
\begin{itemize}[leftmargin=*]
  \item Comprendre le problème de la \emph{redondance sémantique} en
        retrieval pur par similarité.
  \item Connaître la \emph{Maximal Marginal Relevance} (MMR) et ses
        deux paramètres~: \texttt{fetch\_k} et \texttt{lambda\_mult}.
  \item Savoir choisir entre les deux stratégies selon la nature de la
        question (précise vs généraliste).
\end{itemize}
\end{tcolorbox}

\subsection{Solution}

\subsubsection{Architecture}

\begin{center}
\begin{tikzpicture}[node distance=6mm and 10mm]
  \node[data] (q) {"liste vs tuple"};

  \node[block, below left=10mm and 8mm of q] (sim)
       {as\_retriever\\\scriptsize search\_type="similarity"\\\scriptsize k=3};
  \node[block, below right=10mm and 8mm of q] (mmr)
       {as\_retriever\\\scriptsize search\_type="mmr"\\\scriptsize k=3, fetch\_k=20, $\lambda$=0.5};

  \node[doc, below=of sim] (simout) {top 3 par cosinus\\\scriptsize (souvent redondants)};
  \node[doc, below=of mmr] (mmrout) {3 chunks diversifiés\\\scriptsize (parmi 20 candidats)};

  \draw[flowdata] (q)   -- (sim);
  \draw[flowdata] (q)   -- (mmr);
  \draw[flow]     (sim) -- (simout);
  \draw[flow]     (mmr) -- (mmrout);
\end{tikzpicture}
\end{center}

\subsubsection{Le code complet}

\begin{pythoncode}
vectorstore = build_or_load_vectorstore()
question = "Quelle est la différence entre une liste et un tuple ?"

# --- Retriever similarity (défaut) ---
sim_retriever = vectorstore.as_retriever(
    search_type="similarity",
    search_kwargs={"k": 3},
)

# --- Retriever MMR ---
mmr_retriever = vectorstore.as_retriever(
    search_type="mmr",
    search_kwargs={
        "k": 3,                    # nombre final de docs renvoyés
        "fetch_k": 20,             # pool de candidats avant sélection MMR
        "lambda_mult": 0.5,        # 1.0 = pure similarité, 0.0 = pure diversité
    },
)

# --- Comparaison ---
print("--- Top-3 similarity ---")
for i, doc in enumerate(sim_retriever.invoke(question), 1):
    print(f"  [{i}] {doc.page_content[:100]}...")

print("\n--- Top-3 MMR (diversifié) ---")
for i, doc in enumerate(mmr_retriever.invoke(question), 1):
    print(f"  [{i}] {doc.page_content[:100]}...")
\end{pythoncode}

\subsubsection{Explications pas à pas}

\subsubsection*{Étape 1 — Le problème de la similarité pure}

La similarité cosinus classe les chunks par \textbf{proximité avec la
query}, indépendamment les uns des autres. Conséquence~: si la query
matche fortement un certain passage, les \texttt{k} chunks les plus
similaires sont souvent \textbf{très proches entre eux}~:

\begin{quote}\itshape
Pour la query \emph{« liste vs tuple »}, les top-3 par similarité sont
typiquement 3 paragraphes de la \emph{même section} qui parlent de la
même définition formelle. On rate les autres angles du sujet (mutabilité,
performances, conventions d'usage).
\end{quote}

\subsubsection*{Étape 2 — La \emph{Maximal Marginal Relevance}}

MMR sélectionne itérativement les chunks en équilibrant deux critères~:

\begin{itemize}[leftmargin=*]
  \item \textbf{Pertinence}~: similarité avec la query.
  \item \textbf{Diversité}~: \emph{dissimilarité} avec les chunks déjà
        sélectionnés.
\end{itemize}

Algorithme~:
\begin{enumerate}[leftmargin=*]
  \item Récupérer les \texttt{fetch\_k} candidats par similarité pure
        (ex: 20).
  \item Sélectionner le candidat le plus proche de la query (top 1).
  \item Pour chaque candidat restant, calculer~:
        \[
        \text{score} = \lambda \cdot \text{sim}(c, q)
                     - (1-\lambda) \cdot \max_{s \in S} \text{sim}(c, s)
        \]
        où \(S\) est l'ensemble des candidats déjà sélectionnés.
  \item Sélectionner le candidat de plus haut score. Répéter jusqu'à
        \texttt{k}.
\end{enumerate}

\subsubsection*{Étape 3 — Les paramètres \texttt{fetch\_k} et
\texttt{lambda\_mult}}

\begin{center}
\begin{tabular}{p{4cm} p{10cm}}
\toprule
\textbf{Paramètre} & \textbf{Effet} \\
\midrule
\texttt{k} & Nombre final de docs renvoyés. \\[1mm]
\texttt{fetch\_k} & Taille de la pool de candidats sur laquelle MMR
opère. Plus c'est grand, plus la diversité \emph{potentielle} est large.
Typique~: \texttt{fetch\_k = 5 \(\times\) k}. \\[1mm]
\texttt{lambda\_mult} & Curseur similarité $\leftrightarrow$ diversité~:
\texttt{1.0} = MMR redevient similarité pure, \texttt{0.0} = pure
diversité (déconnectée de la query). Sweet spot~: \texttt{0.5} à
\texttt{0.7}. \\
\bottomrule
\end{tabular}
\end{center}

\subsubsection*{Étape 4 — Quand utiliser MMR}

\begin{center}
\begin{tabular}{p{6.5cm} p{6.5cm}}
\toprule
\textbf{\texttt{similarity} (défaut)} & \textbf{\texttt{mmr}} \\
\midrule
Question précise, monolithique
& Question généraliste, multifacette \\[1mm]
« Comment fonctionne \texttt{super()} ? »
& « C'est quoi la POO en Python ? » \\[1mm]
On veut le \emph{meilleur} match exact
& On veut un \emph{tour d'horizon} \\[1mm]
Plus rapide (pas de recalcul)
& Plus lent ($O(k \times \text{fetch\_k})$) \\
\bottomrule
\end{tabular}
\end{center}

\begin{quote}\itshape
\textbf{Règle pratique}~: si tes top-3 par similarité semblent
\emph{redondants} (3 passages disant la même chose), passe en MMR. Si
les top-3 sont déjà variés, similarité suffit (et c'est plus rapide).
\end{quote}

\begin{tcolorbox}[
  colback=white,
  colframe=black!50,
  coltitle=black,
  fonttitle=\footnotesize\bfseries,
  fontupper=\footnotesize,
  title=Sortie de l'exécution,
  breakable,
]
\begin{verbatim}
======================================================================
DEMO 3 -- Similarity vs MMR sur la MEME question
======================================================================

Question : Quelle est la différence entre une liste et un tuple ?

--- Top-3 similarity ---
  [1] , line 1 , in <module> TypeError : 'tuple' object does not
      support item assignment >>> # but they ca...
  [2] >>> list ( x for x in range ( 10 ) if is_even ( x )) [0, 2, 4,
      6, 8] enumerate(iter, start=0) counts...
  [3] elements that are accessed via unpacking (see later in this
      section) or indexing (or even by attribu...

--- Top-3 MMR (diversifié) ---
  [1] , line 1 , in <module> TypeError : 'tuple' object does not
      support item assignment >>> # but they ca...
  [2] 5.1. More on Lists  The list data type has some more methods.
      Here are all of the methods of list o...
  [3] sequences of the same type, the lexicographical comparison is
      carried out recursively.  If all items...

======================================================================
Observe : MMR évite les doublons sémantiques que retourne la
similarity pure. Pour une question généraliste comme 'liste vs
tuple', MMR ramène souvent des chunks plus complémentaires.
======================================================================
\end{verbatim}
\end{tcolorbox}

\subsection{Vérification dans LangSmith}

Pour les deux retrievers, LangSmith affiche le même type de run
(\texttt{VectorStoreRetriever}). La différence se voit dans le
\emph{contenu} des chunks renvoyés~: tu compares visuellement les deux
runs côte-à-côte dans l'UI.

\begin{quote}\itshape
LangSmith n'expose pas (encore) les scores MMR internes. Si tu veux les
inspecter, utilise \texttt{vectorstore.max\_marginal\_relevance\_search\_with\_score(\dots)}
directement, en dehors de la chaîne.
\end{quote}

\subsection*{Récapitulatif}
{\itshape\sloppy
\textit{Ce qu'il faut retenir :}
\begin{enumerate}[leftmargin=*]
  \item La similarité pure peut être \textbf{redondante}~: les top-\texttt{k}
        sont souvent des paragraphes proches du même passage.
  \item MMR équilibre \textbf{pertinence et diversité}~: pioche dans une
        pool de \texttt{fetch\_k} candidats les \texttt{k} les plus
        complémentaires.
  \item \texttt{lambda\_mult} pilote le curseur~: \texttt{0.5}--\texttt{0.7}
        est le sweet spot pour de la doc technique.
  \item Bascule en MMR pour les questions \textbf{généralistes ou
        ouvertes}~; reste en similarity pour les questions précises.
  \item Côté API~: c'est juste un paramètre \texttt{search\_type} à
        changer dans \texttt{.as\_retriever(\dots)}.
\end{enumerate}
\par}

\part{Pour aller plus loin}

% #####################################################################
% #####################################################################
\section{Impact du \texttt{chunk\_size} sur la qualité du RAG}
% #####################################################################
% #####################################################################

\begin{tcolorbox}[title=Exercice bonus 1. Mesurer expérimentalement l'effet du découpage]
Construis trois vectorstores avec \texttt{chunk\_size} = \texttt{500},
\texttt{1000}, \texttt{2000} caractères (et \texttt{chunk\_overlap}
proportionnel, 20\%). Pose la même question précise aux trois et compare
(a) la longueur moyenne des chunks ramenés, (b) leur pertinence visuelle,
(c) la qualité de la réponse finale.

\textbf{Contrainte}~: persister les trois index dans des dossiers
séparés (\texttt{chroma\_db\_500}, etc.) pour éviter de réindexer à
chaque test.

\bigskip
\noindent\textbf{Objectifs pédagogiques}
\begin{itemize}[leftmargin=*]
  \item Voir \textbf{concrètement} pourquoi le \texttt{chunk\_size} est
        l'hyperparamètre n\textsuperscript{o}1 de qualité RAG.
  \item Identifier le compromis \emph{granularité vs contexte}~: petit
        chunk = retrieval précis mais contexte fragmenté ; gros chunk =
        contexte riche mais retriever moins discriminant.
  \item Trouver expérimentalement le sweet spot pour un corpus donné.
\end{itemize}
\end{tcolorbox}

\subsection{Solution}

\subsubsection{Architecture}

\begin{center}
\begin{tikzpicture}[node distance=6mm and 10mm]
  \node[doc] (urls) {URLs Python};
  \node[block, below=of urls] (loader) {WebBaseLoader.load()};

  \node[block, below left=10mm and 14mm of loader] (s500)  {split chunk=500};
  \node[block, below=of loader]                    (s1000) {split chunk=1000};
  \node[block, below right=10mm and 14mm of loader](s2000) {split chunk=2000};

  \node[store, below=of s500]  (db500)  {chroma\_db\_500};
  \node[store, below=of s1000] (db1000) {chroma\_db\_1000};
  \node[store, below=of s2000] (db2000) {chroma\_db\_2000};

  \draw[flow] (urls)   -- (loader);
  \draw[flow] (loader) -- (s500);
  \draw[flow] (loader) -- (s1000);
  \draw[flow] (loader) -- (s2000);
  \draw[flow] (s500)   -- (db500);
  \draw[flow] (s1000)  -- (db1000);
  \draw[flow] (s2000)  -- (db2000);
\end{tikzpicture}
\end{center}

Trois index isolés, même corpus, paramètre de découpage différent. C'est
le protocole expérimental le plus simple pour évaluer un hyperparamètre
en RAG~: \emph{ceteris paribus}, on change une seule variable.

\subsubsection{Le code complet}

\begin{pythoncode}
from pathlib import Path
from langchain_text_splitters import RecursiveCharacterTextSplitter
from langchain_chroma import Chroma
from langchain_community.document_loaders import WebBaseLoader

from phase2_rag import (
    PYTHON_DOCS_URLS, COLLECTION_NAME, get_embeddings, build_rag_chain,
)


def build_vectorstore_with_chunk_size(chunk_size: int, persist_dir: str) -> Chroma:
    """Construit (ou recharge) un vectorstore avec un chunk_size donné."""
    embeddings = get_embeddings()
    if Path(persist_dir).exists():
        return Chroma(
            collection_name=COLLECTION_NAME,
            persist_directory=persist_dir,
            embedding_function=embeddings,
        )

    loader = WebBaseLoader(PYTHON_DOCS_URLS)
    loader.requests_kwargs = {"headers": {"User-Agent": "PyTutor-Edu/0.1"}}
    raw_docs = loader.load()

    splitter = RecursiveCharacterTextSplitter(
        chunk_size=chunk_size,
        chunk_overlap=int(chunk_size * 0.2),       # 20% proportionnel
        separators=["\n\n", "\n", ". ", " ", ""],
    )
    chunks = splitter.split_documents(raw_docs)

    return Chroma.from_documents(
        documents=chunks, embedding=embeddings,
        collection_name=COLLECTION_NAME, persist_directory=persist_dir,
    )


def exercise_1() -> None:
    configs = [
        (500,  "./chroma_db_500"),
        (1000, "./chroma_db_1000"),
        (2000, "./chroma_db_2000"),
    ]
    question = "Comment lève-t-on une exception personnalisée en Python ?"

    for chunk_size, persist_dir in configs:
        print(f"\n--- chunk_size = {chunk_size} ---")
        vs = build_vectorstore_with_chunk_size(chunk_size, persist_dir)
        retriever = vs.as_retriever(search_kwargs={"k": 3})

        docs = retriever.invoke(question)
        avg_len = sum(len(d.page_content) for d in docs) // len(docs)
        print(f"  Longueur moyenne des chunks : {avg_len} caractères")

        chain = build_rag_chain(retriever)
        answer = chain.invoke(question)
        print(f"  Réponse : {answer[:300]}...")
\end{pythoncode}

\subsubsection{Explications pas à pas}

\subsubsection*{Étape 1 — Le compromis granularité vs contexte}

\begin{center}
\begin{tabular}{p{2.5cm} p{5cm} p{5cm}}
\toprule
\textbf{\texttt{chunk\_size}} & \textbf{Avantages} & \textbf{Inconvénients} \\
\midrule
\textbf{500}
& Recherche précise (chaque chunk est ciblé)
& Contexte fragmenté, le modèle manque l'environnement de l'idée \\[1mm]
\textbf{1000} (défaut)
& Bon compromis pour de la doc technique
& — \\[1mm]
\textbf{2000}
& Contexte riche, chunks auto-suffisants
& Beaucoup de bruit, le retriever distingue mal pertinent / tangentiel \\
\bottomrule
\end{tabular}
\end{center}

\subsubsection*{Étape 2 — Mesurer plutôt que deviner}

L'effet du \texttt{chunk\_size} dépend du corpus. Sur la doc Python
(passages bien structurés, paragraphes courts), 1000 est optimal. Sur
des transcripts de conférences (idées qui s'étalent sur plusieurs
paragraphes), 1500--2000 peut être meilleur. \textbf{Mesure} pour ton
corpus.

\begin{quote}\itshape
Le bon outil de mesure n'est pas la longueur ni le nombre de chunks~:
c'est la \textbf{qualité subjective des réponses} sur un échantillon
de questions représentatives. Construis un mini-jeu de 10 questions
typiques et compare les réponses.
\end{quote}

\subsubsection*{Étape 3 — Le \texttt{chunk\_overlap} proportionnel}

On garde \texttt{chunk\_overlap = 0.2 \(\times\) chunk\_size}~: 20\%
de chevauchement entre chunks consécutifs. C'est une bonne valeur
\emph{par défaut} qui évite de couper une idée en deux sans gonfler
excessivement la taille de l'index.

\subsubsection*{Étape 4 — Persister chaque variante}

Réindexer prend du temps. En persistant dans des dossiers séparés
(\texttt{chroma\_db\_500}, etc.), on peut~:
\begin{itemize}[leftmargin=*]
  \item Comparer les variantes rapidement (rechargement $<$1~s chacun).
  \item Garder un environnement reproductible.
  \item Faire le test une fois, l'oublier.
\end{itemize}

\begin{tcolorbox}[
  colback=white,
  colframe=black!50,
  coltitle=black,
  fonttitle=\footnotesize\bfseries,
  fontupper=\footnotesize,
  title=Sortie de l'exécution,
  breakable,
]
\begin{verbatim}
======================================================================
EXERCICE 1 -- Impact du chunk_size sur la qualité du RAG
======================================================================

--- chunk_size = 500 ---
  -> Indexation chunk_size=500 dans ./chroma_db_500...
    -> 658 chunks générés
  Longueur moyenne des chunks retrouvés : 472 caractères
  [1] ) # arguments stored in .args ... print ( inst ) # __str__
      allows args to be pri...
  [2] >>> try : ... raise TypeError ( 'bad type' ) ... except
      Exception as e : ... e ....
  [3] know where to look in case the input came from a file. 8.2.
      Exceptions   Even if...

  Réponse :
  # Lever une exception personnalisée en Python

  D'après le contexte fourni, je peux te montrer comment lever une
  exception en utilisant le mot-clé `raise` [2].

  ## Syntaxe de base

  ```python
  raise TypeException('message')
  ```
  ...

--- chunk_size = 1000 ---
  -> Indexation chunk_size=1000 dans ./chroma_db_1000...
    -> 327 chunks générés
  Longueur moyenne des chunks retrouvés : 981 caractères
  [1] except Exception as inst : ... print ( type ( inst )) # the
      exception type ... p...
  [2] They include SystemExit which is raised by sys.exit() and
      KeyboardInterrupt whic...
  [3] err =} , { type ( err ) =} " ) raise The try ... except
      statement has an optional ...

  Réponse :
  Je dois être honnête : **la réponse à ta question n'est pas
  présente dans le contexte fourni.**

  Le contexte [1], [2] et [3] explique comment :
  - **Capturer** des exceptions avec `try...except`
  - **Accéder** aux informations d'une exception
  - **Re-lever** une exception avec...

--- chunk_size = 2000 ---
  -> Indexation chunk_size=2000 dans ./chroma_db_2000...
    -> 165 chunks générés
  Longueur moyenne des chunks retrouvés : 1977 caractères
  [1] ( inst . args ) # arguments stored in .args ... print ( inst )
      # __str__ allows ...
  [2] # shorthand for 'raise ValueError()' If you need to determine
      whether an excepti...
  [3] +-+---------------- 1 ---------------- | OSError: error 1
      +---------------- 2 --...

  Réponse :
  # Comment lever une exception personnalisée en Python ?

  D'après le contexte fourni, je peux te montrer comment lever une
  exception en Python, mais la documentation ne couvre pas
  explicitement la création d'exceptions personnalisées.
  ...

----------------------------------------------------------------------
OBSERVATIONS A RETENIR :
 - chunk=500  : chunks précis mais souvent incomplets, le modèle
   peut manquer le contexte autour du détail recherché.
 - chunk=1000 : le bon défaut pour de la doc technique.
 - chunk=2000 : beaucoup de bruit, le retriever distingue mal
   les passages pertinents des passages tangentiels.
\end{verbatim}
\end{tcolorbox}

\subsection{Vérification dans LangSmith}

Pour chaque \texttt{chunk\_size}, lance la même question et compare les
runs côte-à-côte. Trois choses se révèlent~:

\begin{enumerate}[leftmargin=*]
  \item \textbf{Tokens du prompt}~: \texttt{chunk\_size=500} avec
        \texttt{k=3} donne $\sim$1500 tokens de contexte ;
        \texttt{chunk\_size=2000} donne $\sim$6000.
  \item \textbf{Pertinence des chunks}~: lisible dans l'output du
        retriever.
  \item \textbf{Qualité de la réponse}~: peut être annotée manuellement
        dans LangSmith (\emph{thumbs up/down}) pour mémoriser l'évaluation.
\end{enumerate}

\begin{quote}\itshape
Le coût en tokens monte vite avec le \texttt{chunk\_size}~: à
\texttt{k=3}, passer de 1000 à 2000 caractères double le coût par
requête sans nécessairement améliorer la qualité. Optimise.
\end{quote}

\subsection*{Récapitulatif}
{\itshape\sloppy
\textit{Ce qu'il faut retenir :}
\begin{enumerate}[leftmargin=*]
  \item Le \texttt{chunk\_size} est l'hyperparamètre n\textsuperscript{o}1
        de qualité RAG. Mesure, ne devine pas.
  \item Sweet spot général pour de la doc technique~: \textbf{800--1200
        caractères}, \texttt{overlap} = 20\%.
  \item Petit chunk~: précis mais fragmenté. Gros chunk~: riche mais
        bruité.
  \item Persiste chaque variante dans son dossier pour comparer sans
        réindexer.
  \item Le coût en tokens augmente linéairement avec \texttt{chunk\_size
        \(\times\) k}.
\end{enumerate}
\par}

% #####################################################################
% #####################################################################
\section{Boucle interactive : transformer le RAG en mini-CLI}
% #####################################################################
% #####################################################################

\begin{tcolorbox}[title=Exercice bonus 2. Une boucle while pour explorer la doc librement]
Transforme la chaîne RAG en une petite REPL (\emph{Read-Eval-Print Loop})~:
une boucle \texttt{while} qui demande une question à l'utilisateur via
\texttt{input()}, appelle la chaîne, affiche la réponse \textbf{et les
sources}, puis recommence. Termine proprement sur ligne vide ou
\texttt{Ctrl+C}.

\textbf{Contrainte}~: utiliser \texttt{build\_rag\_chain\_with\_sources}
de l'exercice 4 pour afficher les sources, et gérer les exceptions
\texttt{EOFError} / \texttt{KeyboardInterrupt} pour une sortie propre.

\bigskip
\noindent\textbf{Objectifs pédagogiques}
\begin{itemize}[leftmargin=*]
  \item Construire une mini-CLI qui sert de banc de test pour explorer
        son corpus librement.
  \item Gérer la terminaison propre (signal et ligne vide).
  \item Préparer mentalement la transition vers une UI web (Gradio,
        Streamlit) où la boucle d'événements est cachée par le framework.
\end{itemize}
\end{tcolorbox}

\subsection{Solution}

\subsubsection{Architecture}

\begin{center}
\begin{tikzpicture}[node distance=6mm and 10mm]
  \node[block] (start) {Démarrage};
  \node[block, below=of start] (load) {build\_or\_load\_vectorstore\\\scriptsize une seule fois};
  \node[block, below=of load] (chain) {build\_rag\_chain\_with\_sources};

  \node[data, below=of chain] (prompt) {\texttt{❯} input()};
  \node[block, below=of prompt] (invoke) {chain.invoke(question)};
  \node[data, below=of invoke] (display) {print(answer)\\print(sources)};

  \draw[flow] (start)  -- (load);
  \draw[flow] (load)   -- (chain);
  \draw[flow] (chain)  -- (prompt);
  \draw[flow] (prompt) -- (invoke);
  \draw[flow] (invoke) -- (display);
  \draw[flow, dashed, draw=orange!70!black] (display.east) ..controls +(2,0) and +(2,0).. (prompt.east)
       node[label, midway, right]{boucle};
\end{tikzpicture}
\end{center}

Trois éléments importants pour la performance perçue~: charger le store
\textbf{une seule fois} (lent), construire la chaîne \textbf{une seule
fois}, et la réutiliser en boucle pour chaque question (rapide).

\subsubsection{Le code complet}

\begin{pythoncode}
from phase2_rag import build_or_load_vectorstore, build_rag_chain_with_sources


def exercise_2_interactive() -> None:
    print("PyTutor — mode interactif (Ctrl+C ou ligne vide pour quitter)")

    # On charge le store et on construit la chaîne UNE SEULE FOIS
    vectorstore = build_or_load_vectorstore()
    retriever   = vectorstore.as_retriever(search_kwargs={"k": 4})
    chain       = build_rag_chain_with_sources(retriever)

    while True:
        try:
            question = input("\n❯ ").strip()
        except (EOFError, KeyboardInterrupt):
            print("\nÀ la prochaine !")
            break

        if not question:
            print("Au revoir.")
            break

        result = chain.invoke(question)

        print("\n" + "─" * 70)
        print(result["answer"])
        print("─" * 70)
        print("Sources :")
        for i, doc in enumerate(result["context"], 1):
            src = doc.metadata.get("source", "?").split("/")[-1]
            print(f"  [{i}] {src}")
\end{pythoncode}

\subsubsection{Explications pas à pas}

\subsubsection*{Étape 1 — Init coûteux, hors de la boucle}

\begin{pythoncode}
vectorstore = build_or_load_vectorstore()         # 1s (rechargement)
retriever   = vectorstore.as_retriever(...)       # instantané
chain       = build_rag_chain_with_sources(...)   # instantané
# Boucle après l'init :
while True:
    ...
\end{pythoncode}

Pattern universel des REPL~: tout ce qui est lent (chargement, parsing,
initialisation) se fait \textbf{une fois} avant la boucle. Dans la
boucle, chaque tour ne fait que ce qui dépend de l'input.

\subsubsection*{Étape 2 — Gérer les sorties propres}

Trois façons pour l'utilisateur de quitter~:

\begin{itemize}[leftmargin=*]
  \item \texttt{Ctrl+C}~: \texttt{input()} lève \texttt{KeyboardInterrupt}.
  \item \texttt{Ctrl+D} (EOF)~: \texttt{input()} lève \texttt{EOFError}.
  \item \textbf{Ligne vide}~: \texttt{input()} renvoie \texttt{""}, qu'on
        traite explicitement.
\end{itemize}

Le \texttt{try/except} attrape les deux premiers, le \texttt{if not
question} attrape le troisième.

\begin{quote}\itshape
Sans \texttt{try/except}, un \texttt{Ctrl+C} propage une trace Python
disgracieuse à l'utilisateur. L'attraper et imprimer un message
amical (« À la prochaine ! ») est la marque du logiciel poli.
\end{quote}

\subsubsection*{Étape 3 — Affichage avec sources}

\texttt{build\_rag\_chain\_with\_sources} renvoie un dict avec
\texttt{answer}, \texttt{context}, \texttt{question}. On affiche~:
\begin{itemize}[leftmargin=*]
  \item La réponse, encadrée d'une ligne pour la séparer visuellement
        du prompt.
  \item Les sources, chacune avec son fichier (\texttt{.split("/")[-1]}
        pour ne garder que le \emph{basename}).
\end{itemize}

\subsubsection*{Étape 4 — Vers une UI web}

Cette mini-CLI est la base mentale d'une UI :
\begin{itemize}[leftmargin=*]
  \item \textbf{Gradio} :
\begin{pythoncode}
import gradio as gr
gr.Interface(
    fn=lambda q: chain.invoke(q)["answer"],
    inputs="text", outputs="text",
).launch()
\end{pythoncode}
  \item \textbf{Streamlit}~: \texttt{st.text\_input}, \texttt{st.write}.
  \item \textbf{FastAPI}~: un endpoint \texttt{POST /ask} qui retourne
        le dict en JSON.
\end{itemize}

Dans tous les cas, la chaîne RAG reste un \texttt{Runnable} inchangé~:
on remplace juste l'interface utilisateur autour.

\begin{tcolorbox}[
  colback=white,
  colframe=black!50,
  coltitle=black,
  fonttitle=\footnotesize\bfseries,
  fontupper=\footnotesize,
  title=Sortie de l'exécution,
  breakable,
]
\begin{verbatim}
======================================================================
EXERCICE 2 -- Mode interactif (Ctrl+C ou ligne vide pour quitter)
======================================================================
[OK] Rechargement de l'index depuis ./chroma_db

PyTutor est prêt. Pose une question sur Python.

> c'est quoi les fonctions décorées?

----------------------------------------------------------------------
# Les fonctions décorées

Une **fonction décorée** est une fonction qui a été transformée par
un **décorateur**. [1]

## Qu'est-ce qu'un décorateur ?

Un décorateur est une fonction qui retourne une autre fonction. Il
est généralement appliqué comme une transformation de fonction en
utilisant la syntaxe `@nom_du_decorateur`. [1]

## Exemple simple

Voici deux façons équivalentes de décorer une fonction : [1]

**Syntaxe avec `@` (sucre syntaxique):**
```python
@staticmethod
def f(arg):
    ...
```

**Syntaxe sans `@` (équivalent):**
```python
def f(arg):
    ...
f = staticmethod(f)
```

Les deux approches font exactement la même chose : elles appliquent
le décorateur `staticmethod` à la fonction `f`.

## Exemples courants

Les décorateurs les plus courants en Python sont : [1]
- `@classmethod`
- `@staticmethod`

Ces décorateurs modifient le comportement de la fonction en la
transformant.
----------------------------------------------------------------------
Sources :
  [1] glossary.html
  [2] errors.html
  [3] controlflow.html
  [4] errors.html

>
Au revoir.
\end{verbatim}
\end{tcolorbox}

\subsection{Vérification dans LangSmith}

Chaque appel à \texttt{chain.invoke(\dots)} produit \textbf{une trace
distincte} dans LangSmith. C'est un super outil pour~:
\begin{itemize}[leftmargin=*]
  \item Repérer les questions qui ont mal marché.
  \item Comparer plusieurs reformulations d'une même question.
  \item Annoter manuellement les bonnes/mauvaises réponses
        (\emph{thumbs up/down}) pour évaluer ton RAG.
\end{itemize}

\subsection*{Récapitulatif}
{\itshape\sloppy
\textit{Ce qu'il faut retenir :}
\begin{enumerate}[leftmargin=*]
  \item Pattern REPL~: init coûteux \textbf{hors} de la boucle, traitement
        rapide \textbf{dans} la boucle.
  \item Gérer la terminaison propre~: \texttt{EOFError},
        \texttt{KeyboardInterrupt}, et ligne vide.
  \item Une CLI interactive est le meilleur banc de test pour explorer
        son corpus et identifier les angles morts.
  \item La chaîne RAG reste inchangée quand on passe à une UI web — c'est
        l'interface qui change, pas le \texttt{Runnable}.
\end{enumerate}
\par}

% #####################################################################
% #####################################################################
\section{Prompt strict pour refuser le hors-contexte}
% #####################################################################
% #####################################################################

\begin{tcolorbox}[title=Exercice bonus 3. Bloquer les hallucinations polies]
Pose une question dont la réponse \textbf{n'est pas} dans les pages
indexées (par exemple~: \emph{« Comment installer un package avec
pip~? »}, sachant que la doc indexée couvre tutoriel + glossaire, pas
l'installation). Observe le comportement par défaut. Puis durcis le
prompt système pour que le modèle refuse explicitement de répondre
hors-contexte.

\textbf{Contrainte}~: garder la même chaîne LCEL, juste swap le prompt.
Le prompt strict doit \textbf{aussi} bien marcher sur des questions
in-context (ne pas être excessivement restrictif).

\bigskip
\noindent\textbf{Objectifs pédagogiques}
\begin{itemize}[leftmargin=*]
  \item Reconnaître l'\emph{hallucination polie}~: le modèle complète
        avec ses connaissances générales quand le contexte ne suffit pas.
  \item Construire un prompt \textbf{défensif} avec règles numérotées
        et exemples de refus.
  \item Vérifier que le refus est exploitable en aval (router vers une
        autre source d'info, par exemple).
\end{itemize}
\end{tcolorbox}

\subsection{Solution}

\subsubsection{Architecture}

\begin{center}
\begin{tikzpicture}[node distance=5mm and 12mm]
  \node[data] (q) {"Comment installer pip ?"};

  \node[block, below left=10mm and 14mm of q]  (def)    {RAG\_PROMPT\\\scriptsize (défaut)};
  \node[block, below right=10mm and 14mm of q] (strict) {STRICT\_RAG\_PROMPT\\\scriptsize (5 règles absolues)};

  \node[llm, below=of def]    (mdef)   {Modèle};
  \node[llm, below=of strict] (mstr)   {Modèle};

  \node[data, below=of mdef, fill=red!10] (rdef)
       {"Pour installer pip, lance\\\textbf{python -m ensurepip\dots}"\\\scriptsize (HALLUCINATION)};
  \node[data, below=of mstr, fill=green!10] (rstr)
       {"Cette information n'est pas\\dans la documentation fournie."\\\scriptsize (REFUS NET)};

  \draw[flowdata] (q)    -- (def);
  \draw[flowdata] (q)    -- (strict);
  \draw[flow]     (def)  -- (mdef);
  \draw[flow]     (strict)-- (mstr);
  \draw[flow]     (mdef) -- (rdef);
  \draw[flow]     (mstr) -- (rstr);
\end{tikzpicture}
\end{center}

\subsubsection{Le code complet}

\begin{pythoncode}
from langchain_core.prompts import ChatPromptTemplate
from langchain_core.output_parsers import StrOutputParser
from langchain_core.runnables import RunnablePassthrough
from config import get_model
from phase2_rag import format_docs


STRICT_RAG_PROMPT = ChatPromptTemplate.from_messages([
    (
        "system",
        "Tu réponds EXCLUSIVEMENT à partir du contexte fourni ci-dessous. "
        "Règles absolues :\n"
        "1. Si la réponse n'apparaît pas littéralement dans le contexte, "
        "tu réponds : 'Cette information n'est pas dans la documentation fournie.' "
        "et tu t'arrêtes.\n"
        "2. Tu ne complètes JAMAIS avec tes connaissances générales sur Python.\n"
        "3. Tu ne fais JAMAIS d'inférence en partant d'un sujet voisin.\n"
        "4. Si tu utilises le contexte, cite les chunks par leur numéro [N].\n"
        "5. Si le contexte est ambigu ou partiel, dis explicitement ce qui "
        "manque plutôt que d'extrapoler."
    ),
    (
        "human",
        "Contexte :\n{context}\n\nQuestion : {question}",
    ),
])


def build_strict_rag_chain(retriever):
    return (
        {
            "context":  retriever | format_docs,
            "question": RunnablePassthrough(),
        }
        | STRICT_RAG_PROMPT
        | get_model()
        | StrOutputParser()
    )
\end{pythoncode}

\subsubsection{Explications pas à pas}

\subsubsection*{Étape 1 — L'hallucination polie}

Sans contraintes explicites, un LLM moderne complète volontiers ses
réponses avec ses connaissances générales~: c'est ce qu'il a appris à
faire pour être « utile ». Sur une question hors-corpus, il va~:

\begin{itemize}[leftmargin=*]
  \item Ignorer (ou paraphraser) le contexte non pertinent.
  \item Sortir une réponse plausible et fluide à partir de ses
        \emph{connaissances pré-entraînées}.
  \item Ne pas mentionner que cette réponse ne vient PAS du contexte.
\end{itemize}

C'est l'\emph{hallucination polie}~: la réponse est techniquement
correcte, mais elle vient d'une source non-traçable. En contexte
pédagogique, c'est exactement ce qu'on veut éviter.

\subsubsection*{Étape 2 — Le pattern \emph{règles absolues numérotées}}

\begin{pythoncode}
"Règles absolues :
1. Si la réponse n'apparaît pas littéralement dans le contexte, ...
2. Tu ne complètes JAMAIS avec tes connaissances générales sur Python.
3. Tu ne fais JAMAIS d'inférence en partant d'un sujet voisin.
4. Si tu utilises le contexte, cite les chunks par leur numéro [N].
5. Si le contexte est ambigu ou partiel, dis explicitement ce qui manque."
\end{pythoncode}

Cinq caractéristiques d'un prompt défensif efficace~:

\begin{itemize}[leftmargin=*]
  \item \textbf{Numérotation}~: les modèles obéissent mieux à des règles
        listées qu'à un paragraphe.
  \item \textbf{Capitales}~: \texttt{EXCLUSIVEMENT}, \texttt{JAMAIS}.
        Effet psychologique mesurable sur les LLM modernes.
  \item \textbf{Phrase exacte de refus}~: « Cette information n'est pas
        dans la documentation fournie. » — le modèle a un \emph{template}
        à recracher, pas à inventer.
  \item \textbf{Couverture explicite des cas voisins}~: règle 3 anti-
        inférence (« cite-moi un sujet proche »), règle 5 anti-vague.
  \item \textbf{Garder la consigne positive}~: règle 4 sur les citations,
        pour ne pas écraser le bon comportement par excès de restrictions.
\end{itemize}

\subsubsection*{Étape 3 — Vérifier la non-restriction excessive}

Un piège classique des prompts stricts~: ils deviennent \emph{trop}
restrictifs et refusent même des questions dont la réponse est dans
le contexte. Test de contrôle indispensable~:

\begin{pythoncode}
in_context = "Qu'est-ce qu'un itérable en Python ?"      # dans le glossaire
print(chain_strict.invoke(in_context))                   # doit répondre normalement
\end{pythoncode}

Si le prompt strict refuse cette question, c'est qu'il est mal calibré
— il faut adoucir le ton sans perdre la rigueur sur le hors-contexte.

\subsubsection*{Étape 4 — Le refus exploitable}

L'intérêt d'un refus net (toujours la même phrase) plutôt que d'un
brouillon hallucinant~: le \emph{code applicatif} peut le détecter
et router~:

\begin{pythoncode}
answer = chain_strict.invoke(question)
if "Cette information n'est pas dans la documentation" in answer:
    # router vers une recherche web, un autre corpus, etc.
    fallback_chain.invoke(question)
\end{pythoncode}

\begin{quote}\itshape
Un bon refus est \textbf{exploitable en aval}. L'hallucination polie ne
l'est pas — il est impossible de savoir programmatiquement si la
réponse est sourcée ou inventée. Toute architecture RAG sérieuse force
le refus.
\end{quote}

\begin{tcolorbox}[
  colback=white,
  colframe=black!50,
  coltitle=black,
  fonttitle=\footnotesize\bfseries,
  fontupper=\footnotesize,
  title=Sortie de l'exécution,
  breakable,
]
\begin{verbatim}
======================================================================
EXERCICE 3 -- Prompt strict vs prompt par défaut
======================================================================
[OK] Rechargement de l'index depuis ./chroma_db

Question (HORS contexte) : Comment installer le package `requests`
                            avec pip ?

--- Prompt PAR DEFAUT ---
Je ne peux pas répondre à cette question en m'appuyant sur le
contexte fourni.

Le contexte contient des informations sur le glossaire Python, le
chemin de recherche des modules (`sys.path`), les objets de type
chemin (`path-like object`), et les packages Python, mais il ne
contient aucune information sur l'installation de packages avec pip.

Pour répondre à ta question sur l'installation de `requests` avec
pip, il faudrait un contexte qui traite de la gestion des packages
et de l'outil pip, ce qui n'est pas disponible ici.

--- Prompt STRICT ---
Cette information n'est pas dans la documentation fournie.

Le contexte fourni traite de concepts comme les modules Python, le
chemin de recherche des modules (`sys.path`), les packages et les
PEPs, mais il ne contient aucune information sur l'installation de
packages avec pip.


Question (DANS contexte, contrôle) : Qu'est-ce qu'un itérable en
                                      Python ?

--- Prompt STRICT ---
Selon la documentation fournie [2], un **itérable** est :

> Un objet capable de retourner ses membres un à un.

Les exemples d'itérables incluent :
- Tous les types de séquences (comme `list`, `str`, et `tuple`)
- Certains types non-séquences comme `dict` et les objets fichier
- Les objets de classes que vous définissez avec une méthode
  `__iter__()` ou une méthode `__getitem__()` qui implémente la
  sémantique de séquence

Les itérables peuvent être utilisés dans une boucle `for` et dans
de nombreux autres contextes où une séquence est nécessaire
(`zip()`, `map()`, etc.). Quand un objet itérable est passé à la
fonction intégrée `iter()`, elle retourne un itérateur pour cet
objet [2].

----------------------------------------------------------------------
OBSERVATIONS A RETENIR :
 - Sans contrainte explicite, le modèle a tendance à 'aider' en
   complétant avec ses connaissances. C'est utile en chat général,
   c'est un BUG en RAG.
 - Le prompt strict rend le refus net et exploitable en aval
   (tu peux router vers une autre source d'info, par exemple).
\end{verbatim}
\end{tcolorbox}

\subsection{Vérification dans LangSmith}

Lance la même question hors-contexte avec les deux prompts. Dans
LangSmith~:
\begin{enumerate}[leftmargin=*]
  \item Compare les \emph{outputs} côte-à-côte~: l'un parle de pip,
        l'autre refuse.
  \item Compare les \emph{system prompts}~: le strict a 4-5 fois plus
        de tokens (pour mémoire, $\sim$200 tokens supplémentaires).
  \item Compare les \emph{tokens de sortie}~: le refus est court ($\sim$30
        tokens), l'hallucination est longue ($\sim$300 tokens). Le prompt
        strict réduit la facture sur les requêtes hors-corpus.
\end{enumerate}

\subsection*{Récapitulatif}
{\itshape\sloppy
\textit{Ce qu'il faut retenir :}
\begin{enumerate}[leftmargin=*]
  \item L'\emph{hallucination polie} est le piège n\textsuperscript{o}1
        du RAG~: le modèle complète sans le dire avec ses connaissances
        générales.
  \item Prompt défensif efficace~: règles \textbf{numérotées},
        \textbf{capitales} sur les interdictions, phrase exacte de refus,
        couverture des cas voisins.
  \item Toujours tester avec une question \textbf{in-context} aussi, pour
        ne pas créer un prompt trop restrictif.
  \item Un refus net (phrase canonique) est \textbf{exploitable en aval}~:
        on peut router vers une autre source.
  \item Coût~: prompt système $\sim$200 tokens en plus, mais réponse
        plus courte sur hors-contexte $\to$ facture similaire ou inférieure.
\end{enumerate}
\par}

% #####################################################################
% #####################################################################
\section{Vérifier et exploiter LangSmith}
% #####################################################################
% #####################################################################

\begin{tcolorbox}[title=Exercice bonus 4. Activer le tracing et lire l'arbre d'exécution]
Configure LangSmith dans ton \texttt{.env}, exécute une chaîne RAG, et
explore l'arbre d'exécution dans l'UI \texttt{smith.langchain.com}.
Vérifie programmatiquement la présence des variables d'environnement
nécessaires.

\textbf{Contrainte}~: l'exercice doit afficher un message d'aide clair
si la config est manquante (variables \texttt{LANGSMITH\_TRACING},
\texttt{LANGSMITH\_API\_KEY}, \texttt{LANGSMITH\_PROJECT}), avec les
lignes à ajouter au \texttt{.env}.

\bigskip
\noindent\textbf{Objectifs pédagogiques}
\begin{itemize}[leftmargin=*]
  \item Configurer LangSmith~: variables d'environnement et endpoint
        EU vs US.
  \item Lire l'arbre d'exécution d'une chaîne RAG~: les branches du
        \texttt{RunnableParallel}, le \texttt{format\_docs}, le prompt
        rendu, l'appel modèle, le parser.
  \item Apprendre à \textbf{utiliser LangSmith comme principale source
        de vérité} sur le comportement d'une chaîne (pas le code, pas
        les logs).
\end{itemize}
\end{tcolorbox}

\subsection{Solution}

\subsubsection{Architecture}

\begin{center}
\begin{tikzpicture}[node distance=5mm and 14mm]
  \node[data] (env) {.env\\\scriptsize LANGSMITH\_TRACING=true\\\scriptsize LANGSMITH\_API\_KEY=...\\\scriptsize LANGSMITH\_PROJECT=PyTutor2};
  \node[block, below=of env] (script) {python script.py};
  \node[block, below=of script] (chain) {chain.invoke(\dots)};

  \node[block, right=20mm of chain] (api) {API LangSmith};
  \node[store, above=of api] (cloud) {smith.langchain.com\\\scriptsize traces persistées};

  \draw[flowdata] (env)    -- (script);
  \draw[flow]     (script) -- (chain);
  \draw[flow]     (chain)  -- (api) node[label, midway, above]{trace envoyée};
  \draw[flow]     (api)    -- (cloud);
\end{tikzpicture}
\end{center}

\subsubsection{Le code complet}

\begin{pythoncode}
import os
from phase2_rag import build_or_load_vectorstore, build_rag_chain_with_sources


def exercise_4_langsmith() -> None:
    # --- Vérification de la configuration ---
    tracing = os.getenv("LANGSMITH_TRACING", "").lower() in {"true", "1"}
    api_key = os.getenv("LANGSMITH_API_KEY")
    project = os.getenv("LANGSMITH_PROJECT", "default")

    print(f"LANGSMITH_TRACING : {'✓ activé' if tracing else '✗ NON activé'}")
    print(f"LANGSMITH_API_KEY : {'✓ présente' if api_key else '✗ absente'}")
    print(f"LANGSMITH_PROJECT : {project}")

    if not (tracing and api_key):
        print("→ Ajoute dans .env :")
        print("    LANGSMITH_TRACING=true")
        print("    LANGSMITH_API_KEY=lsv2_pt_...")
        print("    LANGSMITH_PROJECT=PyTutor2")
        print("    LANGSMITH_ENDPOINT=https://eu.api.smith.langchain.com  # si compte EU")
        return

    # --- Exécution d'une chaîne pour générer une trace ---
    vectorstore = build_or_load_vectorstore()
    retriever   = vectorstore.as_retriever(search_kwargs={"k": 4})
    chain       = build_rag_chain_with_sources(retriever)

    result = chain.invoke("Comment écrire un générateur infini en Python ?")
    print(result["answer"][:300] + "...")
\end{pythoncode}

\subsubsection{Explications pas à pas}

\subsubsection*{Étape 1 — Les 3 variables d'environnement essentielles}

\begin{center}
\begin{tabular}{p{4.5cm} p{10cm}}
\toprule
\textbf{Variable} & \textbf{Rôle} \\
\midrule
\texttt{LANGSMITH\_TRACING}
& \texttt{"true"} pour activer l'envoi automatique des traces. Sans ça,
  les chaînes s'exécutent normalement mais rien n'est envoyé. \\[1mm]
\texttt{LANGSMITH\_API\_KEY}
& Clé personnelle (commence par \texttt{lsv2\_pt\_\dots}). À créer sur
  \texttt{smith.langchain.com} (gratuit). \\[1mm]
\texttt{LANGSMITH\_PROJECT}
& Nom du projet où les traces seront groupées. Ex.~: \texttt{PyTutor2}. \\[1mm]
\texttt{LANGSMITH\_ENDPOINT} & Optionnel.
\texttt{https://eu.api.smith.langchain.com} pour un compte EU
(défaut~: US). \\
\bottomrule
\end{tabular}
\end{center}

\subsubsection*{Étape 2 — Vérification programmatique}

Lire \texttt{os.getenv} avec un défaut explicite, et tester~:

\begin{pythoncode}
tracing = os.getenv("LANGSMITH_TRACING", "").lower() in {"true", "1"}
api_key = os.getenv("LANGSMITH_API_KEY")
\end{pythoncode}

Avantage~: l'utilisateur voit \emph{exactement} ce qui manque et a les
lignes à coller dans son \texttt{.env}.

\subsubsection*{Étape 3 — L'arbre dans l'UI}

Sur \texttt{smith.langchain.com}, ouvre ton projet et clique sur ta
trace la plus récente. Le panneau gauche montre l'arbre d'exécution
avec, pour la chaîne RAG~:

\begin{enumerate}[leftmargin=*]
  \item Le \texttt{RunnableParallel} (les 2 branches \texttt{context}
        et \texttt{question}).
  \item Le \texttt{Retriever} (avec sa query exacte et la liste de docs).
  \item Le \texttt{format\_docs} (un \texttt{RunnableLambda}).
  \item Le \texttt{ChatPromptTemplate} (avec le prompt rendu, texte
        complet envoyé au LLM).
  \item L'\texttt{appel modèle} (avec tokens consommés, latence, coût \$).
  \item Le \texttt{StrOutputParser}.
\end{enumerate}

\subsubsection*{Étape 4 — Pourquoi LangSmith est plus utile que les
logs Python}

\begin{center}
\begin{tabular}{p{6.5cm} p{6.5cm}}
\toprule
\textbf{Logs Python (\texttt{print}, \texttt{logging})} & \textbf{LangSmith} \\
\midrule
Manuel~: tu choisis quoi logger.
& Automatique~: tracé pour chaque \texttt{Runnable}. \\[1mm]
Texte brut, dispersé sur stdout/fichier.
& UI structurée, recherchable, partagable. \\[1mm]
Pas de versionnage ni d'historique.
& Toutes les traces archivées avec timestamp. \\[1mm]
Pas de coût/latence agrégés.
& Onglets dédiés \emph{Costs}, \emph{Latency}, \emph{Tokens}. \\[1mm]
Tu vois ce que TU as choisi de logger.
& Tu vois \textbf{exactement ce que le modèle a vu}. \\
\bottomrule
\end{tabular}
\end{center}

\begin{quote}\itshape
LangSmith devient la principale source de vérité dès qu'on travaille
sur une chaîne sérieuse. Les \texttt{print} restent utiles pour le
local quick-and-dirty, mais la moindre investigation passe par l'UI.
\end{quote}

\begin{tcolorbox}[
  colback=white,
  colframe=black!50,
  coltitle=black,
  fonttitle=\footnotesize\bfseries,
  fontupper=\footnotesize,
  title=Sortie de l'exécution,
  breakable,
]
\begin{verbatim}
======================================================================
EXERCICE 4 -- Tracer une exécution dans LangSmith
======================================================================

LANGSMITH_TRACING : [OK] activé
LANGSMITH_API_KEY : [OK] présente
LANGSMITH_PROJECT : PyTutor2

Exécution d'une chaîne RAG (la trace sera envoyée à LangSmith)...
[OK] Rechargement de l'index depuis ./chroma_db

Réponse :
# Comment écrire un générateur infini en Python ?

Je dois te dire que **la réponse à cette question n'est pas
présente dans le contexte fourni**.

Le contexte explique comment fonctionnent les générateurs en
général [3,4] : ce sont des fonctions qui utilisent `yield` pour
retourner des valeurs une...

----------------------------------------------------------------------
CE QU'IL FAUT REGARDER DANS L'UI :
  1. Va sur https://smith.langchain.com et ouvre le projet 'PyTutor2'
  2. Tu vois ta dernière trace. Clique dessus.
  3. Le panneau de gauche montre l'arbre d'exécution :
     - le RunnableParallel 'context' + 'question' (les 2 branches
       partent en parallèle)
     - puis le format_docs (RunnableLambda)
     - puis le prompt formatté (regarde le texte exact envoyé)
     - puis l'appel au modèle (avec coût et latence)
     - puis le parser
  4. Clique sur l'appel modèle pour voir les tokens consommés
     et le prompt exact, sans rien deviner.
\end{verbatim}
\end{tcolorbox}

\subsection{Vérification dans LangSmith}

\emph{(Cet exercice EST la vérification.)}

À regarder dans l'UI~:
\begin{itemize}[leftmargin=*]
  \item L'onglet \emph{Runs} liste toutes les traces récentes.
  \item Le panneau de \emph{détail} d'une trace montre l'arbre.
  \item L'onglet \emph{Costs} montre le coût \$ et les tokens.
  \item L'onglet \emph{Annotations} permet de noter manuellement
        les bonnes/mauvaises réponses (\emph{thumbs up/down}).
\end{itemize}

\subsection*{Récapitulatif}
{\itshape\sloppy
\textit{Ce qu'il faut retenir :}
\begin{enumerate}[leftmargin=*]
  \item 3 variables d'env essentielles~: \texttt{LANGSMITH\_TRACING=true},
        \texttt{LANGSMITH\_API\_KEY=lsv2\_pt\_\dots},
        \texttt{LANGSMITH\_PROJECT=ton-projet}.
  \item Compte EU~: ajouter
        \texttt{LANGSMITH\_ENDPOINT=https://eu.api.smith.langchain.com}.
  \item Vérification programmatique au démarrage de l'app~: meilleure
        UX que de découvrir l'absence de trace plus tard.
  \item L'UI LangSmith devient la principale source de vérité~: arbre
        d'exécution, prompts rendus, coût, latence.
  \item Annoter manuellement avec \emph{thumbs up/down} crée un dataset
        d'évaluation gratuit.
\end{enumerate}
\par}

% #####################################################################
% #####################################################################
\section{\texttt{MultiQueryRetriever} : reformuler la question}
% #####################################################################
% #####################################################################

\begin{tcolorbox}[title=Exercice bonus 5. Améliorer le recall avec un LLM en amont]
Remplace le \texttt{Retriever} simple par un \texttt{MultiQueryRetriever}
qui reformule chaque question utilisateur en 3 variantes via un LLM
\textbf{avant} la recherche, puis fusionne les résultats des 3 recherches.
Compare sur une question volontairement courte/ambiguë (ex.~: \emph{«
Comment marche l'héritage~? »}).

\textbf{Contrainte}~: activer le logging du module
\texttt{langchain.retrievers.multi\_query} pour voir les reformulations
générées par le LLM dans le terminal.

\bigskip
\noindent\textbf{Objectifs pédagogiques}
\begin{itemize}[leftmargin=*]
  \item Comprendre la notion de \emph{recall} en information retrieval
        (retrouver toutes les bonnes infos, pas juste les plus
        pertinentes en surface).
  \item Utiliser un LLM pour \textbf{augmenter} la question utilisateur
        (et non pour répondre).
  \item Mesurer le coût supplémentaire (1 appel LLM de reformulation +
        N recherches au lieu d'une).
  \item Savoir quand utiliser MultiQuery~: questions courtes / ambiguës ;
        et quand ne PAS l'utiliser~: questions précises et bien formulées.
\end{itemize}
\end{tcolorbox}

\subsection{Solution}

\subsubsection{Architecture}

\begin{center}
\begin{tikzpicture}[node distance=4mm and 10mm]
  \node[data] (q) {"Comment marche l'héritage ?"};

  \node[llm, below=of q] (gen) {LLM\\\scriptsize génère 3 reformulations};

  \node[data, below=of gen, font=\scriptsize] (vars)
       {"Comment hériter en Python ?"\\
        "Syntaxe de la classe enfant en Python"\\
        "Différence entre super() et héritage"};

  \node[block, below left=10mm and 8mm of vars]  (r1) {Retriever};
  \node[block, below=of vars]                    (r2) {Retriever};
  \node[block, below right=10mm and 8mm of vars] (r3) {Retriever};

  \node[doc, below=of r1] (d1) {docs};
  \node[doc, below=of r2] (d2) {docs};
  \node[doc, below=of r3] (d3) {docs};

  \node[block, below=of d2, minimum width=72mm] (merge) {Union + déduplication};
  \node[doc, below=of merge] (final) {chunks uniques};

  \draw[flowdata] (q)    -- (gen);
  \draw[flow]     (gen)  -- (vars);
  \draw[flow]     (vars) -- (r1);
  \draw[flow]     (vars) -- (r2);
  \draw[flow]     (vars) -- (r3);
  \draw[flow]     (r1)   -- (d1);
  \draw[flow]     (r2)   -- (d2);
  \draw[flow]     (r3)   -- (d3);
  \draw[flow]     (d1)   -- (merge);
  \draw[flow]     (d2)   -- (merge);
  \draw[flow]     (d3)   -- (merge);
  \draw[flow]     (merge)-- (final);
\end{tikzpicture}
\end{center}

Pipeline en quatre temps~: reformulation LLM, N recherches parallèles,
union, déduplication.

\subsubsection{Le code complet}

\begin{pythoncode}
import logging
from langchain_classic.retrievers.multi_query import MultiQueryRetriever
from config import get_model
from phase2_rag import build_or_load_vectorstore, build_rag_chain


def exercise_5() -> None:
    vectorstore     = build_or_load_vectorstore()
    base_retriever  = vectorstore.as_retriever(search_kwargs={"k": 3})
    multi_retriever = MultiQueryRetriever.from_llm(
        retriever=base_retriever,
        llm=get_model(),
    )

    # Question courte et ambiguë : "héritage" peut couvrir plusieurs angles
    question = "Comment marche l'héritage ?"

    # --- Active le logging pour voir les reformulations ---
    logging.basicConfig()
    logging.getLogger("langchain.retrievers.multi_query").setLevel(logging.INFO)

    # --- Retriever simple ---
    docs_simple = base_retriever.invoke(question)
    print(f"--- Simple (k=3) ---")
    for i, d in enumerate(docs_simple, 1):
        print(f"  [{i}] {d.page_content[:100]}...")

    # --- MultiQuery (les logs montrent les 3 reformulations) ---
    docs_multi = multi_retriever.invoke(question)
    print(f"\n--- MultiQuery ({len(docs_multi)} docs après déduplication) ---")
    for i, d in enumerate(docs_multi, 1):
        print(f"  [{i}] {d.page_content[:100]}...")

    # --- Comparaison des réponses finales ---
    print("\n--- RÉPONSE simple ---")
    print(build_rag_chain(base_retriever).invoke(question)[:400], "...")

    print("\n--- RÉPONSE MultiQuery ---")
    print(build_rag_chain(multi_retriever).invoke(question)[:400], "...")
\end{pythoncode}

\subsubsection{Explications pas à pas}

\subsubsection*{Étape 1 — \emph{Precision} vs \emph{Recall}}

\begin{center}
\begin{tabular}{p{4cm} p{10cm}}
\toprule
\textbf{Métrique} & \textbf{Définition (en RAG)} \\
\midrule
\textbf{Precision}
& Parmi les chunks retournés, combien sont pertinents~? \\[1mm]
\textbf{Recall}
& Parmi les chunks pertinents qui existent dans le corpus, combien
sont retrouvés~? \\
\bottomrule
\end{tabular}
\end{center}

Le retriever simple optimise la \emph{precision} (top-k par similarité
exacte). \texttt{MultiQuery} optimise le \emph{recall}~: en multipliant
les formulations de la question, on \textbf{couvre plus d'angles} et on
retrouve des chunks qui étaient invisibles à la formulation originale.

\subsubsection*{Étape 2 — Comment la reformulation marche}

\texttt{MultiQueryRetriever.from\_llm} envoie un prompt au LLM~:

\begin{quote}\itshape
« Génère 3 reformulations différentes de la question suivante pour
améliorer une recherche par similarité~: "Comment marche l'héritage~?" »
\end{quote}

Le LLM produit (par exemple)~:
\begin{itemize}[leftmargin=*]
  \item « Comment hériter d'une classe en Python~? »
  \item « Quelle est la syntaxe d'une classe enfant qui hérite~? »
  \item « Différence entre héritage et composition en Python »
\end{itemize}

Chaque variante est ensuite passée au retriever de base. Les résultats
sont \textbf{fusionnés et dédupliqués} (par identité de chunk).

\subsubsection*{Étape 3 — Le logging interne}

\begin{pythoncode}
import logging
logging.basicConfig()
logging.getLogger("langchain.retrievers.multi_query").setLevel(logging.INFO)
\end{pythoncode}

Avec ces 3 lignes, les reformulations générées s'affichent dans la
console à chaque appel. Indispensable pour~:
\begin{itemize}[leftmargin=*]
  \item Comprendre ce que le LLM \emph{reformule}.
  \item Diagnostiquer si les reformulations sont pertinentes ou
        triviales.
  \item Ajuster le \texttt{prompt} de reformulation au besoin.
\end{itemize}

\subsubsection*{Étape 4 — Coût et quand l'utiliser}

\begin{center}
\begin{tabular}{p{5cm} p{4cm} p{5cm}}
\toprule
\textbf{Métrique} & \textbf{Simple} & \textbf{MultiQuery} \\
\midrule
Appels LLM pour la query
& 0 (embedding seul)
& 1 (reformulation) \\[1mm]
Appels au retriever
& 1
& \texttt{N} (typique~: 3) \\[1mm]
Coût par requête utilisateur
& \(\sim\)1\texttt{x}
& \(\sim\)1.5--2\texttt{x} \\[1mm]
Gain sur questions précises
& —
& \textbf{Nul}, juste du coût \\[1mm]
Gain sur questions courtes/ambiguës
& —
& \textbf{Significatif} (+30--50\% de recall) \\
\bottomrule
\end{tabular}
\end{center}

\begin{quote}\itshape
\textbf{Règle pratique}~: utilise MultiQuery quand tu observes que des
réponses sont incomplètes \emph{parce que des chunks pertinents n'ont
pas été ramenés}. N'utilise pas MultiQuery si le problème est ailleurs
(chunk\_size, prompt, modèle).
\end{quote}

\begin{tcolorbox}[
  colback=white,
  colframe=black!50,
  coltitle=black,
  fonttitle=\footnotesize\bfseries,
  fontupper=\footnotesize,
  title=Sortie de l'exécution,
  breakable,
]
\begin{verbatim}
[OK] Rechargement de l'index depuis ./chroma_db

Question (volontairement vague) : Comment marche l'héritage ?

--- Retriever SIMPLE (k=3) ---
  [1] 'UnicodeError', 'UnicodeTranslateError', 'UnicodeWarning',
      'UserWarning', 'ValueError', 'Warning', '...
  [2] |     raise ExceptionGroup('there were problems', excs) |
      ExceptionGroup: there were problems (2 sub...
  [3] definition.  (A class is never used as a global scope.)  While
      one rarely encounters a good reason f...

--- MultiQuery Retriever ---

  11 documents au total (déduplication automatique)
  [1] 4.9. More on Defining Functions 4.9.1. Default Argument Values
      4.9.2. Keyword Arguments 4.9.3. Speci...
  [2] Previous topic Enum HOWTO Next topic Logging HOWTO This page
      Report a bug Improve this page Show sou...
  [3] mixed numeric types are compared according to their numeric
      value, so 0 equals 0.0, etc.  Otherwise,...
  [4] . text encoding   A string in Python is a sequence of Unicode
      code points (in range U+0000 - U+10FFF...
  [5] Encoding and error handler used by Python to decode bytes from
      the operating system and encode Unico...
  [6] simply replace the base class method of the same name. There is
      a simple way to call the base class ...
  [7] There's nothing special about instantiation of derived classes:
      DerivedClassName() creates a new ins...
  [8] |     raise ExceptionGroup('there were problems', excs) |
      ExceptionGroup: there were problems (2 sub...
  [9] cycles so that the memory can be reclaimed. data race   A
      situation where multiple threads access th...
  [10] yield and a return statement is that on reaching a yield the
       generator's state of execution is suspe...
  [11] __iter__() method which returns an object with a __next__()
       method.  If the class defines __next__()...

--- REPONSE avec retriever simple ---
# L'héritage en Python

D'après le contexte fourni [3], voici comment fonctionne l'héritage :

## Syntaxe de base

Pour créer une classe qui hérite d'une autre, on utilise la syntaxe
suivante :

```python
class DerivedClassName(BaseClassName):
    <statement-1>
    ...
    <statement-N>
```

**Explication :**
- `DerivedClassName` est la classe dérivée (la classe enfant)
- `BaseClassName` est la cl...

--- REPONSE avec MultiQuery ---
# L'héritage en Python

D'après la documentation fournie [6][7], voici comment fonctionne
l'héritage :

## Syntaxe de base

Une classe dérivée hérite d'une ou plusieurs classes de base :

```python
class ClasseDerivee(ClasseBase):
    <instruction-1>
    ...
    <instruction-N>
```

## Fonctionnement principal

**Résolution des méthodes** [7] : Quand vous appelez une méthode,
Python cherche l'attr...

----------------------------------------------------------------------
OBSERVATIONS A RETENIR :
 - Regarde les logs ci-dessus pour voir les 3 reformulations
   générées automatiquement par le LLM.
 - MultiQuery double ou triple ton coût par question (1 appel pour
   reformuler + N appels de recherche), mais améliore notablement
   le recall sur les questions courtes ou floues.
 - A ne PAS utiliser sur des questions très précises et bien
   formulées : aucun gain, juste du coût en plus.
\end{verbatim}
\end{tcolorbox}

\subsection{Vérification dans LangSmith}

Avec MultiQuery, LangSmith trace un arbre plus profond~:

\begin{center}
\begin{tikzpicture}[every node/.style={font=\footnotesize\ttfamily, anchor=west}]
  \node                              (a) {MultiQueryRetriever};
  \node[below=1mm of a, xshift=4mm]  (b) {LLMChain (génération des 3 reformulations) ← LLM \#1};
  \node[below=1mm of b]              (c) {VectorStoreRetriever (variante 1)};
  \node[below=1mm of c]              (d) {VectorStoreRetriever (variante 2)};
  \node[below=1mm of d]              (e) {VectorStoreRetriever (variante 3)};
  \node[below=1mm of e, xshift=-4mm] (f) {(fusion + déduplication, hors trace)};
\end{tikzpicture}
\end{center}

\begin{quote}\itshape
Tu vois directement les 3 variantes générées (onglet \emph{Output} de
la \texttt{LLMChain}), et le contenu des 3 retrievals. Précieux pour
auditer la qualité des reformulations.
\end{quote}

\subsection*{Récapitulatif}
{\itshape\sloppy
\textit{Ce qu'il faut retenir :}
\begin{enumerate}[leftmargin=*]
  \item \texttt{MultiQueryRetriever} reformule la question en N variantes
        \textbf{via un LLM}, puis fusionne les N résultats du retriever
        de base.
  \item Objectif~: améliorer le \emph{recall} sur les questions courtes,
        ambiguës ou mal formulées.
  \item Coût~: +1 appel LLM pour reformuler + N appels au retriever.
        $\sim$1.5 à 2$\times$ le coût d'une requête simple.
  \item Active le \texttt{logging} sur \texttt{langchain.retrievers.multi\_query}
        pour voir les reformulations.
  \item À ne PAS utiliser sur des questions déjà précises~: aucun gain,
        juste du coût en plus.
\end{enumerate}
\par}

% #####################################################################
% #####################################################################
\section{Filtrage par métadonnée}
% #####################################################################
% #####################################################################

\begin{tcolorbox}[title=Exercice bonus 6. Restreindre la recherche à une sous-partie du corpus]
Chaque chunk du \texttt{Chroma} a hérité d'une \texttt{metadata.source}
(l'URL d'origine). Construis un retriever filtré qui ne cherche que dans
\texttt{tutorial/classes.html} pour une question sur la POO. Compare
avec le retriever non filtré.

Démontre aussi le \textbf{piège}~: si tu poses une question hors du
périmètre du filtre (ex.~: « Comment lever une exception~? » alors que
le filtre est sur \texttt{classes.html}), le retriever ramène du contenu
hors-sujet.

\textbf{Contrainte}~: utiliser \texttt{search\_kwargs=\{"filter": \{\dots\}\}}
avec la syntaxe MongoDB-like de Chroma.

\bigskip
\noindent\textbf{Objectifs pédagogiques}
\begin{itemize}[leftmargin=*]
  \item Utiliser la \texttt{metadata} des chunks pour restreindre la
        recherche.
  \item Connaître la syntaxe de filtre Chroma (\texttt{\$eq},
        \texttt{\$in}, \texttt{\$and}, \texttt{\$or}, \dots).
  \item Identifier le piège~: un filtre mal calibré force le retriever
        à choisir le « moins pire » dans un sous-corpus inadapté.
  \item Esquisser le \emph{pattern de production}~: un classifieur LLM
        en amont qui décide \emph{dynamiquement} du filtre à appliquer.
\end{itemize}
\end{tcolorbox}

\subsection{Solution}

\subsubsection{Architecture}

\begin{center}
\begin{tikzpicture}[node distance=6mm and 10mm]
  \node[data] (q) {"Comment fonctionne super() ?"};

  \node[block, below=of q] (ret) {Retriever\\\scriptsize filter=\{source: classes.html\}};

  \node[store, below=of ret] (chroma)
       {Chroma\\\scriptsize tous les chunks};

  \node[doc, right=of chroma, fill=red!8] (filtered) {chunks de\\classes.html\\\scriptsize uniquement};

  \draw[flowdata] (q)     -- (ret);
  \draw[flow]     (ret)   -- (chroma);
  \draw[flow]     (chroma)-- (filtered) node[label, midway, above]{filtre metadata};
\end{tikzpicture}
\end{center}

Le filtre s'applique \emph{avant} le calcul de similarité~: seuls les
chunks dont les metadata matchent sont éligibles à la recherche.

\subsubsection{Le code complet}

\begin{pythoncode}
from phase2_rag import build_or_load_vectorstore


def exercise_6() -> None:
    vectorstore = build_or_load_vectorstore()

    # --- Retriever filtré : SEULEMENT classes.html ---
    classes_only = vectorstore.as_retriever(
        search_kwargs={
            "k": 4,
            "filter": {"source": "https://docs.python.org/3/tutorial/classes.html"},
        },
    )

    # --- Retriever non filtré pour comparaison ---
    unfiltered = vectorstore.as_retriever(search_kwargs={"k": 4})

    question = "Comment fonctionne `super()` ?"

    print("--- SANS filtre ---")
    for i, d in enumerate(unfiltered.invoke(question), 1):
        src = d.metadata.get("source", "?").split("/")[-1]
        print(f"  [{i}] ({src}) {d.page_content[:80]}...")

    print("\n--- AVEC filtre source=classes.html ---")
    for i, d in enumerate(classes_only.invoke(question), 1):
        src = d.metadata.get("source", "?").split("/")[-1]
        print(f"  [{i}] ({src}) {d.page_content[:80]}...")

    # Démo de la limite : question hors du périmètre du filtre
    off_topic = "Comment lever une exception ?"
    print(f"\n--- AVEC filtre, question hors-périmètre ({off_topic}) ---")
    for i, d in enumerate(classes_only.invoke(off_topic), 1):
        src = d.metadata.get("source", "?").split("/")[-1]
        print(f"  [{i}] ({src}) {d.page_content[:80]}...")
\end{pythoncode}

\subsubsection{Explications pas à pas}

\subsubsection*{Étape 1 — La syntaxe Chroma}

\begin{pythoncode}
search_kwargs={
    "k": 4,
    "filter": {"source": "https://docs.python.org/3/tutorial/classes.html"},
}
\end{pythoncode}

Forme simple~: \texttt{\{key: value\}} = match exact sur la metadata.
Sous le capot, Chroma traduit en \texttt{\{key: \{\$eq: value\}\}}.

Pour des filtres plus riches, syntaxe MongoDB-like~:

\begin{center}
\begin{tabular}{ll}
\toprule
\textbf{Opérateur} & \textbf{Exemple} \\
\midrule
\texttt{\$eq}, \texttt{\$ne}
& \texttt{\{"source": \{"\$eq": "url"\}\}} \\
\texttt{\$in}, \texttt{\$nin}
& \texttt{\{"source": \{"\$in": [url1, url2]\}\}} \\
\texttt{\$gt}, \texttt{\$gte}, \texttt{\$lt}, \texttt{\$lte}
& sur des champs numériques \\
\texttt{\$and}, \texttt{\$or}
& \texttt{\{"\$and": [filtre1, filtre2]\}} \\
\bottomrule
\end{tabular}
\end{center}

\subsubsection*{Étape 2 — Quand le filtre brille}

Cas d'usage parfaits~:
\begin{itemize}[leftmargin=*]
  \item \textbf{Multi-tenant}~: chaque utilisateur a son propre corpus,
        filtrer par \texttt{tenant\_id}.
  \item \textbf{Multi-langue}~: filtrer par \texttt{lang} pour éviter
        de retrouver du contenu en anglais sur une query française.
  \item \textbf{Documentation versionnée}~: filtrer par \texttt{version}
        (« super() en Python 3.0 vs 3.12 »).
  \item \textbf{Dates}~: récupérer uniquement les chunks plus récents
        que \texttt{X}.
\end{itemize}

\subsubsection*{Étape 3 — Quand le filtre casse}

Si le filtre est \textbf{mal calibré} (utilisateur cherche dans le
mauvais sous-corpus), le retriever ramène du contenu hors-sujet sans
warning~:

\begin{quote}\itshape
Sur \emph{« Comment lever une exception~? »} avec filter
\texttt{classes.html}, Chroma retourne les chunks de \texttt{classes.html}
les plus proches de la query — qui parlent de classes, pas d'exceptions.
Le RAG aval va générer une réponse hors-sujet ou hallucinée à partir
de ce contexte inadapté.
\end{quote}

C'est le danger du filtre~: il \textbf{force} le retrieval à choisir
dans un sous-ensemble, même si rien n'y est pertinent.

\subsubsection*{Étape 4 — Pattern de production : auto-routing}

En production, le filtre ne devrait pas être codé en dur. Pattern
typique~:

\begin{center}
\begin{tikzpicture}[node distance=5mm and 12mm]
  \node[data] (q) {question utilisateur};
  \node[llm, below=of q] (cls) {Classifieur LLM\\\scriptsize choisit la catégorie};
  \node[data, below=of cls] (cat) {\{"category": "exceptions"\}};
  \node[block, below=of cat] (filter) {build filter};
  \node[block, below=of filter] (ret) {Retriever filtré\\\scriptsize ad hoc};

  \draw[flowdata] (q)    -- (cls);
  \draw[flow]     (cls)  -- (cat);
  \draw[flow]     (cat)  -- (filter);
  \draw[flow]     (filter) -- (ret);
\end{tikzpicture}
\end{center}

Le mini-classifieur LLM (avec \texttt{with\_structured\_output} +
\texttt{Literal[\dots]} comme dans la Phase 1 / Exercice 4 bonus) décide
\emph{dynamiquement} du filtre à appliquer. C'est l'analogue du
\texttt{RunnableBranch} mais appliqué au retrieval.

\begin{tcolorbox}[
  colback=white,
  colframe=black!50,
  coltitle=black,
  fonttitle=\footnotesize\bfseries,
  fontupper=\footnotesize,
  title=Sortie de l'exécution,
  breakable,
]
\begin{verbatim}
======================================================================
EXERCICE 6 -- Filtrage par métadonnée (source URL)
======================================================================
[OK] Rechargement de l'index depuis ./chroma_db

Question : Comment fonctionne `super()` ?

--- SANS filtre (toutes les pages indexées) ---
  [1] (modules.html) 'UnicodeError', 'UnicodeTranslateError',
      'UnicodeWarning', 'UserWarning', 'Value...
  [2] (modules.html) 'EOFError', 'Ellipsis', 'EnvironmentError',
      'Exception', 'False', 'FileExistsErr...
  [3] (modules.html) Function   The built-in function dir() is used
      to find out which names a module ...
  [4] (glossary.html) . single dispatch   A form of generic function
      dispatch where the implementation...

--- AVEC filtre source=classes.html ---
  [1] (classes.html) definition.  (A class is never used as a global
      scope.)  While one rarely encoun...
  [2] (classes.html) call-next-method and is more powerful than the
      super call found in single-inheri...
  [3] (classes.html) __iter__() method which returns an object with
      a __next__() method.  If the clas...
  [4] (classes.html) # Function defined outside the class def f1 (
      self , x , y ): return min ( x , x...


Démo limite -- question hors-périmètre : Comment lever une exception ?

--- AVEC filtre source=classes.html (mauvaise idée ici) ---
  [1] (classes.html) 285 >>> xvec = [ 10 , 20 , 30 ] >>> yvec = [ 7 ,
      5 , 3 ] >>> sum ( x * y for x ,...
  [2] (classes.html) 9. Classes 9.1. A Word About Names and Objects
      9.2. Python Scopes and Namespaces...
  [3] (classes.html) Examples, recipes, and other code in the
      documentation are additionally licensed...
  [4] (classes.html) __iter__() method which returns an object with
      a __next__() method.  If the clas...

----------------------------------------------------------------------
OBSERVATIONS A RETENIR :
 - Le filtre par metadata est un couteau : utile quand tu SAIS
   où chercher (ex: namespace, tag, langue, date), pénalisant
   quand tu te trompes de domaine.
 - Syntaxe Chroma : opérateurs $eq, $ne, $in, $gt, $gte, $lt,
   $lte, $and, $or. Ex: {'source': {'$in': [url1, url2]}}.
 - Pattern de prod : un mini-classifieur LLM en amont qui décide
   du filtre à appliquer (auto-routing par metadata).
\end{verbatim}
\end{tcolorbox}

\subsection{Vérification dans LangSmith}

Le run du retriever filtré expose le \texttt{filter} dans son input~:

\begin{center}
\begin{tikzpicture}[every node/.style={font=\footnotesize\ttfamily, anchor=west}]
  \node                              (a) {VectorStoreRetriever (filtered)};
  \node[below=1mm of a, xshift=4mm]  (b) {Input  : "Comment fonctionne super() ?"};
  \node[below=1mm of b]              (c) {Config : k=4, filter=\{source: ".../classes.html"\}};
  \node[below=1mm of c]              (d) {Output : [Doc, Doc, Doc, Doc]};
  \node[below=1mm of d, xshift=4mm]  (e) {tous source: classes.html};
\end{tikzpicture}
\end{center}

\begin{quote}\itshape
Vérifie toujours dans LangSmith que la \texttt{metadata.source} des
chunks retournés est bien dans le périmètre attendu. Si tu vois des
sources hors du filtre, c'est qu'il y a un bug dans la spécification
(souvent une typo dans l'URL).
\end{quote}

\subsection*{Récapitulatif}
{\itshape\sloppy
\textit{Ce qu'il faut retenir :}
\begin{enumerate}[leftmargin=*]
  \item \texttt{search\_kwargs=\{"filter": \{\dots\}\}} restreint le
        retrieval aux chunks dont la \texttt{metadata} matche.
  \item Syntaxe Chroma MongoDB-like~: \texttt{\$eq}, \texttt{\$in},
        \texttt{\$and}, etc. Match exact si on omet l'opérateur.
  \item Bons cas d'usage~: multi-tenant, multi-langue, versionnage,
        dates.
  \item Piège~: si la question sort du périmètre du filtre, le retriever
        ramène du contenu hors-sujet sans warning. Toujours vérifier.
  \item Pattern de production~: classifieur LLM en amont qui décide
        \emph{dynamiquement} du filtre à appliquer (auto-routing par
        metadata).
\end{enumerate}
\par}



\appendix

\section{Programme principal exécutable}

La fonction \texttt{demo()} ci-dessous orchestre les trois démos
présentées tout au long du chapitre~: indexation (ou rechargement),
RAG simple, RAG avec sources, puis comparaison \texttt{similarity} vs
\texttt{mmr}. C'est le point d'entrée exécuté par
\texttt{python phase2\_rag.py}.

\begin{pythoncode}
def demo() -> None:
    # --- Indexation (ou rechargement) ---
    vectorstore = build_or_load_vectorstore()

    # --- Démo 1 : RAG simple, sortie texte ---
    print("=" * 70)
    print("DÉMO 1 — RAG simple (retriever similarity, k=4)")
    print("=" * 70)
    retriever = vectorstore.as_retriever(
        search_type="similarity",
        search_kwargs={"k": 4},
    )
    chain = build_rag_chain(retriever)

    question = DEMO_QUESTIONS[1]
    print(f"\nQuestion : {question}\n")
    answer = chain.invoke(question)
    print(answer)

    # --- Démo 2 : RAG avec sources ---
    print("\n" + "=" * 70)
    print("DÉMO 2 — RAG avec inspection des sources")
    print("=" * 70)
    chain_with_sources = build_rag_chain_with_sources(retriever)
    question = DEMO_QUESTIONS[2]
    print(f"\nQuestion : {question}\n")
    result = chain_with_sources.invoke(question)
    print("RÉPONSE :")
    print(result["answer"])
    print("\nSOURCES UTILISÉES :")
    for i, doc in enumerate(result["context"], 1):
        snippet = doc.page_content[:120].replace("\n", " ")
        print(f"  [{i}] {doc.metadata.get('source', '?')}")
        print(f"      {snippet}...")

    # --- Démo 3 : comparaison similarity vs MMR ---
    print("\n" + "=" * 70)
    print("DÉMO 3 — Similarity vs MMR sur la MÊME question")
    print("=" * 70)
    question = DEMO_QUESTIONS[0]
    print(f"\nQuestion : {question}")

    sim_retriever = vectorstore.as_retriever(
        search_type="similarity",
        search_kwargs={"k": 3},
    )
    mmr_retriever = vectorstore.as_retriever(
        search_type="mmr",
        # fetch_k : on récupère 20 candidats avant de sélectionner
        # les 3 plus diversifiés.
        # lambda_mult : 1.0 = pure similarité, 0.0 = pure diversité.
        search_kwargs={"k": 3, "fetch_k": 20, "lambda_mult": 0.5},
    )

    print("\n--- Top-3 similarity ---")
    for i, doc in enumerate(sim_retriever.invoke(question), 1):
        print(f"  [{i}] {doc.page_content[:100].replace(chr(10), ' ')}...")

    print("\n--- Top-3 MMR (diversifié) ---")
    for i, doc in enumerate(mmr_retriever.invoke(question), 1):
        print(f"  [{i}] {doc.page_content[:100].replace(chr(10), ' ')}...")

    print("\n" + "=" * 70)
    print("Observe : MMR évite les doublons sémantiques que retourne la")
    print("similarity pure. Pour une question généraliste comme 'liste vs")
    print("tuple', MMR ramène souvent des chunks plus complémentaires.")
    print("=" * 70)

if __name__ == "__main__":
    demo()
\end{pythoncode}

\end{document}
